\documentclass[11pt]{article}

\usepackage[margin=1in]{geometry}
\usepackage{amsmath,amssymb,amsthm,mathtools}
\usepackage{enumitem}
\usepackage{booktabs}
\usepackage{tabularx,longtable,array}
\usepackage{microtype}
\usepackage[hidelinks]{hyperref}
\usepackage{xurl}
\usepackage[nameinlink,noabbrev]{cleveref}

% Build-hygiene: the worked-example wrapper is identifier-heavy.  Use ragged X
% columns and modest emergency stretch so long JSON/route ids wrap instead of
% producing thousands of overfull boxes in the Hold-lane compile sweep.
\newcolumntype{Y}{>{\raggedright\arraybackslash}X}
\setlength{\emergencystretch}{3em}


\newcommand{\MaxL}{\mathcal{L}}

\newcommand{\TV}{\mathrm{TV}}

\newcommand{\PP}{\mathbb{P}}

\title{Worked Example: Instantiating a Tiered Reader-Privacy Receipt\\
\large Receipt line items, user discovery, and evaluation interlocks}
\author{}
% NOTE: The shipped worked-example artifacts (receipt, TW/ENF/state registries, validation report)
% correspond to note_version=1.99 and release_id=worked-example-draft-199. This paper is an
% explanatory wrapper around that bundle and should not drift from the artifact ids.
\date{Unified archive note v2.09 (2026-06-18; archive rev0900; worked-example-draft-199)}


\begin{document}
\maketitle

\begin{abstract}
This note is an end-to-end audit specimen for one reader-privacy claim: a relayed delegated-routing primary path with a budgeted structured-overlay fallback.
It contributes no new anonymity theorem.
Its purpose is to show how a terse public paper can name a claim, expose a few replayable line items, and stop at the first sufficient owner without publishing the full artifact bundle inline. The maintained attenuation slice is arithmetic-only: its receipt explicitly says that no joint observation-channel witness is present and that the slice is not publication-eligible.

The worked cut is intentionally small in public view.
A referee can inspect six stops---receipt, owner/root map, replay plan, release spine, answer-side disclosure packet, and review menu---then widen only if the question genuinely changes.
The supporting archive remains requestable, and the wrapper now treats the paper version separately from the maintained artifact cut: the shipped objects are still \texttt{worked-example-draft-199} / \texttt{note\_version=1.99}, while this explanatory wrapper may advance to repair prose or validation checks without changing those private witnesses.
\end{abstract}

\paragraph{How to read this note.}
Skip everything except:
(i) \Cref{sec:scenario} for the declared threat window and exposure surfaces,
(ii) \Cref{sec:lineitems} for the six representative receipt line items and their replay hooks,
(iii) \Cref{tab:workedexample-auditwalk} for the first-sufficient audit stops,
(iv) \Cref{sec:drill} for one concrete audience-question walk, and
(v) \Cref{sec:numerical} for the tiny arithmetic slice that downstream tools consume.
All other sections are there only to show the interfaces compose without ``semantic drift'' and without owner-boundary drift.

\paragraph{Validation boundary.}
The maintained witness cut and the paper wrapper deliberately have different version clocks. The artifacts remain pinned at \texttt{worked-example-draft-199} / \texttt{note\_version=1.99}; the wrapper version records explanatory and release-spine repairs. The local validator should therefore check that the wrapper names the pinned \texttt{release\_id} and artifact \texttt{note\_version}, not that the wrapper's prose version equals the artifact note version. It also checks that the maintained receipt digest is coherent across the UVI, ABOM, support manifest, and release receipt, that active worked-route references no longer cite an older receipt digest, that Eval~2 active advantage-carrier rows still name the maintained receipt digest, release id, artifact note version, exposure normal form, detector recipe, TV budget, and quotable transfer interval, and that the Eval~1 / Eval~1A notary and support carriers remain aligned with the structured verifier packet and checkpoint-support tuple.

\paragraph{Owned background (avoid redundancy).}
Receipt tuples and JSONPath meanings are owned by Synthesis~12; ENF IDs and observation tiers by Synthesis~3 and Synthesis~15; threat windows by Synthesis~4; primary-path separation guards by Synthesis~22--23; replay-plan minimality and plan/state equivalence by Synthesis~18--20; user discovery by Synthesis~13; and endpoint conversion by Synthesis~5 together with the posterior-mass/PML endpoint (Anonymity~A).\cite{series:receiptitems,series:traceifaces,series:tieredobs,series:threatwindows,series:primarysepmanifest,series:probsep,series:planstateequiv,series:auditorreplay,series:userminiface,series:endpointbridge,series:oddsinflation}
Artifact availability and disclosure are governed by ABOM/OINL and the series artifact policy; the exact paper-facing source-only packet/menu cuts, completion criteria, public-status tokens, carryforward rules, outward-facing notices, canonical notice families, token-keyed notice-family selections, reusable response packets, packet carryforward rules, packet delta ledgers, refresh notices, refresh-notice normal forms, refresh-notice selections, successor follow-up menus, successor-facing refresh-response packets, and refreshed-packet closure verdicts remain owned by their respective source-only successor notes rather than by this worked paper.\cite{series:abomoinl,series:artifactpolicy,series:publicrequestcontracts,series:requestableevidenceclasses,series:requestfulfillmentcertificates,series:publicrequeststatusenvelopes,series:publicrequestcarryforward,series:publicrequestnotices,series:publicrequestnoticenormalforms,series:publicrequestnoticeselections,series:publicrequestresponsepackets,series:publicrequestresponsepacketcarryforward,series:publicrequestresponsepacketdeltas,series:publicrequestresponsepacketrefreshnotices,series:publicrequestresponsepacketrefreshnoticenormalforms,series:publicrequestresponsepacketrefreshnoticeselections,series:publicrequestresponsepacketrefreshresponsemenus,series:publicrequestresponsepacketrefreshresponsepackets,series:publicrequestresponsepacketrefreshresponsepacketclosureverdicts}
For the five published mathematical backbones and their canonical wiki links, use Synthesis~24 rather than re-explaining theorem roots inline. Stable release-evolution stage ownership belongs to Synthesis~32--34, the theorem-to-review bridge that turns those local stages into one stable series spine belongs to Synthesis~55, and the wider theorem/worked/release/notary/claim/release stop ladder belongs to Synthesis~76. This note imports those neighbors so the worked artifacts can stay concrete rather than encyclopedic.\cite{series:mathcrosswalk,series:releasespine,series:releaseclosure,series:statusenvelopes,series:seriesspines,series:theoremtopublicationspine}

\paragraph{Later-terse-paper citation posture.}
Later terse papers should cite this worked note only when the live issue is one concrete auditable route through the maintained bundle or one maintained worked cut inside that route: a bridge-cut card, a packet-local ladder, the answer-side terminal form, the worked publication-boundary card, the first bounded omitted-support packet just beyond the archive cut, one exact request-side rung, or the request-side terminal form. They should instead cite Synthesis~74 when the live issue is the abstract paired answer/request sentence discipline, Synthesis~75 when the live issue is the abstract stop/widen/reopen judgment for that tail, Synthesis~24 for theorem-root routing, Synthesis~32 for release-local stage order, Synthesis~55 for the abstract series spine, and Synthesis~76 when the live issue has widened into theorem/worked/release/notary/claim/release stop selection rather than one concrete maintained route. That is the only way this note stays an example rather than becoming a shadow owner for the whole tail.\cite{series:mathcrosswalk,series:releasespine,series:seriesspines,series:pairedterminalcitationdiscipline,series:pairedterminalcutoverrules,series:theoremtopublicationspine}

\paragraph{Bridge-cut worked-surface reuse.}
Before a later terse paper widens to any smaller packet-local, answer-side, or source-only worked cut, it may first stop at one maintained bridge-cut card in \nolinkurl{artifacts/example_series_spine.json}. The older \texttt{minimal\_bridge\_cuts} rows now carry their own \texttt{bridge\_cut\_sentence\_normal\_form}, \texttt{bridge\_cut\_reopen\_artifact}, and \texttt{bridge\_cut\_reopen\_rule}, so the worked data itself now says how a terse paper may quote the theorem/receipt/replay/release/archive-cut bridge and where widening must reopen next. Later terse papers should therefore stop at the first sufficient maintained bridge-cut card before they widen to packet-local ladders, answer-side terminal cuts, adjacent owner buckets, exact request rungs, or the full audience/question route.\cite{series:seriesspines}

\paragraph{Line-item-facing worked-surface reuse.}
The maintained worked bundle now exposes more than one reusable concrete cut. For emitted receipt or verifier objects, \nolinkurl{artifacts/example_series_spine.json} exports packet-local ladders under \texttt{line\_item\_spines.*}, and later terse papers should use those ladders only when the live issue is one exact carried packet together with its immediate theorem/replay/release neighbors rather than the whole audience/question route. Those packet-local ladders now carry not only \texttt{stop\_when}, \texttt{widen\_only\_if}, and \texttt{route\_local\_reopen\_artifact}, but also a quotable \texttt{packet\_local\_sentence\_normal\_form} and an explicit \texttt{route\_local\_reopen\_rule}, so the worked data itself now says both how a terse paper may quote the packet-local cut and exactly how that shortcut fails closed once the live issue has widened past the line-item bridge. For owner-map-only rows, the same list is deliberately a boundary cut rather than a packet-local carrier: rows such as \texttt{line\_item\_spines.retry\_schedule\_packet}, \texttt{line\_item\_spines.runtime\_conformance\_packet}, and \texttt{line\_item\_spines.routing\_signature\_manifest} now carry \texttt{materialization\_status=owner\_map\_spine\_only\_current\_cut}, \texttt{replay\_plan\_id\_status=route\_target\_only\_not\_plan\_catalog\_row\_current\_cut}, and \texttt{public\_objects\_status=route\_context\_only\_not\_packet\_materialization\_evidence}. Later terse papers may cite those rows to confirm the route split, not to claim emitted receipt, UVI, verifier, replay-plan, or support-bundle packets.\cite{series:operationalA,series:operationalB,series:certifiedB,series:seriesspines}

\paragraph{Base receipt owner-map support rows.}
The carried base receipt rows now label their owner-map support boundary as explicitly as the profiling row does. In \nolinkurl{example_line_item_owner_map.json}, \nolinkurl{primary_path_declaration}, \nolinkurl{fallback_contact_surface}, \nolinkurl{fallback_tiered_vector}, \nolinkurl{path_selection_indicator}, and \nolinkurl{profiling_equalization} carry \nolinkurl{materialization_status=base_receipt_current_cut}, \nolinkurl{receipt_file=example_receipt.json}, a replay-plan id, \nolinkurl{plan_spec_id}, a support-bundle id, \nolinkurl{support_consumed_by}, \nolinkurl{support_required_artifacts}, \nolinkurl{support_pointers}, and packet/replay/support public objects. The support-bundle ids are now explicit: \nolinkurl{support_bundle_id=bundle://worked-example/primary}, \nolinkurl{support_bundle_id=bundle://worked-example/fallback-cct}, \nolinkurl{support_bundle_id=bundle://worked-example/fallback-vector}, \nolinkurl{support_bundle_id=bundle://worked-example/pathselect}, and \nolinkurl{support_bundle_id=bundle://worked-example/profiling-eq}. The support-bundle resolver is now multiplicity-safe: each materialized bundle row carries \nolinkurl{resolved_owner_map_row_names} and \nolinkurl{resolved_owner_map_rows}, and \nolinkurl{bundle://worked-example/fallback-cct} resolves both \nolinkurl{fallback_contact_surface} and the verifier-carried \nolinkurl{notary_replay_packet} without singleton \nolinkurl{line_item_name}/\nolinkurl{materialization_status} shortcut fields. The generated \nolinkurl{line_item_spines} rows mirror the replay \nolinkurl{plan_spec_id} and \nolinkurl{support_anchor} payloads, so the owner-map row, receipt line item, replay plan, and support-bundle map now have the same current-cut boundary without requiring a reviewer to infer the support payload from two separate files. The verifier report now closes the same base-plan surface: \nolinkurl{example_verifier_report.json} lists all five base receipt replay hooks in \nolinkurl{plans_checked}, including \nolinkurl{eval-profiling-eq-v1}, and \nolinkurl{support_artifacts.profiling_evidence_digest=f9b06968e2985a425a8e05830b31ead21ba2715c6567330e0f2500677798187e}. In this cut the \nolinkurl{plans_checked rows now mirror support_bundle_id}, \nolinkurl{support_consumed_by}, \nolinkurl{support_required_artifacts}, \nolinkurl{support_pointers}, \nolinkurl{receipt_file}, \nolinkurl{public_objects}, and \nolinkurl{public_object_digest_entries} from the matching materialized owner-map rows. The \nolinkurl{public_object_digest_entries} field binds every named verifier public object to its current in-archive SHA-256 digest, so a verifier-plan citation is no longer a thin plan-id-only list or a filename-only support list. The \nolinkurl{verifier_contract} field now also binds \nolinkurl{canonical_input_order}, \nolinkurl{public_object_digest_entries}, \nolinkurl{evaluator_source_sha256}, \nolinkurl{materializer_source_sha256}, and \nolinkurl{proof_carrying_status=formal_exact_rational_budget_closure_log_endpoint_derivations_and_profiling_cp_union_bound_present}; this status is deliberately scoped: the maintained exact-rational closure certificate is formal for the public decimal budget endpoints, and the companion transcendental endpoint derivation bundle proves the maintained \texttt{log2} endpoint upper bounds by exact rational series and integer witnesses; the profiling eta/delta finite-sample row is now certified by the exact-rational CP/union-bound bundle. The optional profiling variant verifier surface stays separate from base \nolinkurl{plans_checked}: \nolinkurl{example_verifier_report.json} now carries \nolinkurl{optional_variant_plans_checked} for \nolinkurl{eval-profiling-prefixfetch-v1}, \nolinkurl{eval-profiling-prefixfetch-statecond-v1}, and \nolinkurl{eval-profiling-contact-coarsened-v1}; each row mirrors \nolinkurl{receipt_variant_file}, \nolinkurl{evidence_id}, \nolinkurl{support_bundle_id}, \nolinkurl{support_required_artifacts}, \nolinkurl{support_pointers}, \nolinkurl{public_objects}, \nolinkurl{public_object_digest_entries}, and \nolinkurl{verifier_contract} from the matching variant owner-map row. For every materialized owner-map row, the declared \nolinkurl{support_required_artifacts are a subset of public_objects}; this is now a validator-checked closure rule rather than a reader convention. The verifier-carried \nolinkurl{notary_replay_packet} is labeled separately with \nolinkurl{materialization_status=verifier_report_current_cut}, \nolinkurl{verifier_report_file=example_verifier_report.json}, \nolinkurl{replay_plan_id=eval-fallback-cct-v1}, \nolinkurl{plan_spec_id=sha256:e9168dbec2cabd9115a7f3ea9aebe0a75ff910d2295e450a340716cc7e541193}, \nolinkurl{support_bundle_id=bundle://worked-example/fallback-cct}, and the matching \nolinkurl{support_required_artifacts}/\nolinkurl{support_pointers}; its \nolinkurl{public_objects} list now also includes \nolinkurl{example_receipt.json} and \nolinkurl{example_state_decl.json}, so it should not be mistaken for a base receipt line item even though it shares the fallback contact replay route.

\paragraph{Answer-side terminal worked-surface reuse.}
Before a later terse paper ever crosses the archive cut, it should stop at the answer-side terminal form itself rather than reopen the whole audience/question route. The maintained \texttt{answer\_spines} exported in \nolinkurl{artifacts/example\_series\_spine.json} and mirrored in \nolinkurl{artifacts/example\_question\_routes.json} now carry not only the four-anchor \texttt{answer\_side\_terminal\_citation\_normal\_form} and its one-clause \texttt{answer\_side\_terminal\_sentence\_normal\_form}, but also an explicit \texttt{answer\_side\_terminal\_stop\_when}, an \texttt{answer\_side\_terminal\_widen\_only\_if}, an \texttt{answer\_side\_terminal\_reopen\_artifact}, and an \texttt{answer\_side\_terminal\_first\_reopen\_owner\_map}. So the worked data itself now says when the answer-side shortcut is still honest, when a paper may widen into the adjacent source-only branch, and which imported answer-side owner must be reopened first if exact-location, full-audit, or reveal-order becomes live. The larger one is route-local: \nolinkurl{artifacts/example_question_routes.json} and the neighboring paper-facing route objects remain the honest stop only when the live issue has become the exact checked-answer card, the first bounded omitted-support packet, or the stop/widen/reopen judgment for one maintained audience question.

\paragraph{Adjacent owner-bucket worked-surface reuse.}
Once a later terse paper has honestly left the answer-side archive cut, it still should not jump directly to the request-side capstone. The maintained worked spine mirrors four adjacent source-only owner buckets under \texttt{minimal\_bridge\_cuts} and \texttt{source\_only\_owner\_buckets}: \texttt{paper\_level\_request\_contracts\_and\_notice\_choice}, \texttt{notice\_keyed\_response\_packetization}, \texttt{visible\_refresh\_notice\_choice}, and \texttt{successor\_refresh\_response\_packetization\_and\_closure}. Later terse papers should stop at the first such bucket whose \texttt{stop\_when} already matches the live question, render that stop by the bucket's own \texttt{bucket\_sentence\_normal\_form}, and widen only by that bucket's \texttt{widen\_only\_if}. If the bucket label is still too coarse, the same worked object now says exactly where to reopen next through \texttt{route\_local\_reopen\_artifact}, \texttt{route\_local\_first\_rung\_owner}, and \texttt{route\_local\_first\_rung\_artifact}.\cite{series:seriesspines}

\paragraph{Exact source-only ladder-rung reuse.}
Once that first sufficient bucket is fixed, later terse papers should still stop again at the first exact source-only request-ladder rung whose \texttt{stop\_when} already matches the live question. The maintained \texttt{source\_only\_request\_ladder} entries in \nolinkurl{artifacts/example\_series\_spine.json} now carry their own \texttt{rung\_sentence\_normal\_form} together with \texttt{route\_local\_reopen\_artifact} and \texttt{route\_local\_reopen\_rule}, so the worked data itself now says both how a terse paper may quote one exact rung and where it must go next if the generic rung label is still too coarse for the instantiated audience/question route. The paired request-side terminal shortcut is therefore a late stop inside the fourth bucket, not permission to skip the earlier bucket choices or the exact rung-level stop that make the widened source-only branch auditable in the first place.\cite{series:seriesspines}

\paragraph{Paired-terminal cutover worked-surface reuse.}
Later terse papers should also treat the stop-versus-widen-versus-reopen judgment itself as its own smaller worked cut before reopening the whole instantiated route. The maintained \texttt{answer\_spines} in \nolinkurl{artifacts/example\_series\_spine.json} and the mirrored audience/question routes in \nolinkurl{artifacts/example\_question\_routes.json} now carry a \texttt{paired\_terminal\_cutover\_sentence\_normal\_form} together with \texttt{paired\_terminal\_stop\_when}, \texttt{paired\_terminal\_request\_terminal\_artifact}, \texttt{paired\_terminal\_request\_terminal\_kind}, \texttt{paired\_terminal\_request\_terminal\_widen\_only\_if}, \texttt{paired\_terminal\_reopen\_artifact}, and \texttt{paired\_terminal\_fail\_closed\_reopen\_rule}. So the worked data itself now says when a later note may still quote the cutover card, what request-side terminal floor that card points at, when widening is honest only after the first sufficient adjacent owner bucket and exact request rung are already fixed, and where fail-closed reopening must go when one named trigger fires. The full audience/question route remains the honest stop only when even that cutover card is no longer enough because the instantiated reopened owner rather than the cutover judgment itself has become the live issue.\cite{series:pairedterminalcutoverrules,series:workedexample}

\paragraph{A minimal public worked-example audit-stop card.}
\begin{table}[t]
\centering
\small
\begin{tabularx}{\linewidth}{@{}Y Y Y@{}}
\toprule
Live question & Smallest quotable stop & Widen only if \\
\midrule
Which compact public answer is actually being shipped for this compared release pair? & Answer-side terminal form in \nolinkurl{example_series_spine.json}; quote its own \texttt{answer\_side\_terminal\_sentence\_normal\_form} or \texttt{answer\_side\_terminal\_citation\_normal\_form} for \nolinkurl{example_compare_report.json} $\to$ \texttt{auditor\_compact\_diff\_summary\_answer\_card} $\to$ \texttt{auditor\_compact\_diff\_summary\_answer\_disclosure\_packet} $\to$ \texttt{auditor\_compact\_diff\_summary\_answer\_review\_menu}. & The answer spine's own \texttt{answer\_side\_terminal\_widen\_only\_if} or \texttt{answer\_side\_terminal\_reopen\_rule} becomes live: exact location, full audit trail, reveal order, or a question that has already crossed the answer-side archive cut. \\
Which adjacent source-only owner now answers once the question has crossed the archive cut? & First adjacent owner bucket in \nolinkurl{example_series_spine.json}; stop at the first of \texttt{paper\_level\_request\_contracts\_and\_notice\_choice} $\to$ \texttt{notice\_keyed\_response\_packetization} $\to$ \texttt{visible\_refresh\_notice\_choice} $\to$ \texttt{successor\_refresh\_response\_packetization\_and\_closure} whose \texttt{stop\_when} already matches, then quote its \texttt{bucket\_sentence\_normal\_form}; if one exact rung is already fixed inside that bucket, stop again at the first matching \texttt{source\_only\_request\_ladder} entry and quote its \texttt{rung\_sentence\_normal\_form}. & One follow-up class, one exact bounded packet, one exact closure slice, or a need to reopen the first route-local rung becomes the live issue. \\
What is the first honest source-only packet once one follow-up class is fixed? & First bounded omitted-support packet: \nolinkurl{example_public_request_notice.json} $\to$ \texttt{worked\_example\_receipt\_interlock\_successor\_public\_request\_response\_menu\_\_request\_scope\_check}. & The live issue becomes successor-visible refresh closure rather than the first bounded omitted-support slice. \\
Where does the current visible-refresh branch close? & Request-side terminal form: opening notice $\to$ first bounded omitted-support packet $\to$ exact successor-facing packet $\to$ exact closure verdict. & A named reopen trigger fires, so terminal shortcuts are no longer honest for this branch. \\
When is the paired stop/widen/reopen judgment itself the live issue? & Paired-terminal cutover card in \nolinkurl{example_series_spine.json}; quote its own \texttt{paired\_terminal\_cutover\_sentence\_normal\_form} together with the request-side floor named by \texttt{paired\_terminal\_request\_terminal\_artifact} / \texttt{paired\_terminal\_request\_terminal\_kind}. & The cutover card's own \texttt{paired\_terminal\_request\_terminal\_widen\_only\_if} or \texttt{paired\_terminal\_fail\_closed\_reopen\_rule} becomes live because the first sufficient bucket / exact rung is not yet fixed or a named reopen trigger has already fired. \\
\bottomrule
\end{tabularx}
\caption{A minimal public worked-example audit-stop card. This is the smallest public answer later terse papers can quote directly when the live issue is one concrete auditable route through the maintained example rather than the theorem root, generic release spine, or wider source-only maintenance tail.}
\label{tab:workedexample-mincard}
\end{table}

\paragraph{Tiny cutover drill.}
Suppose a later short note only needs to justify the public sentence \emph{``the replay-required successor stayed on the same certified interface and the shipped compact-diff answer is the right public summary.''} The honest citation packet is the answer-side terminal form in \Cref{tab:workedexample-mincard}, not the whole six-stop walk. If the referee then asks \emph{which source-only owner answers once the question has crossed the archive cut?}, stop first at the earliest adjacent owner bucket whose \texttt{stop\_when} already matches, quote that bucket's own sentence form, and only reopen the first route-local rung named by that same bucket if the bucket label is still too coarse. If the next live issue has already become one exact request-contract / status / notice / packet / closure rung inside that bucket, stop again at the first matching \texttt{source\_only\_request\_ladder} entry, quote that rung's own sentence form, and only then reopen the audience/question route if the rung label itself is still too coarse, rather than jumping straight to the request-side capstone. Only after the follow-up class is fixed as \texttt{request\_scope\_check} should the paper widen one step further to the first bounded request-side packet. Only when the live question becomes \emph{where does the visible-refresh branch close?} should the paper widen to the request-side terminal form; if one named reopen trigger fires, stop using the shortcut and reopen the named owner instead.

\section{Scenario and declared interfaces}
\label{sec:scenario}
We instantiate one deployment slice with two paths:
\begin{enumerate}[leftmargin=*]
\item \textbf{Primary: relayed delegated routing.} A relay separates client-side metadata from destination-side metadata (a primary separation claim guarded by a published failure probability $\rho$). Deployed control-plane surfaces for delegated routing (AutoConf, overrides, path scopes) are pinned by the state declaration per Synthesis~30.\cite{series:primarysepmanifest,series:deployedsurfaces}
\item \textbf{Fallback: structured overlay lookup.} When primary fails (or is bypassed by policy), a structured overlay lookup is executed and may expose a Tier~2 committee-contact trace (CCT) together with coarse timing and congestion witnesses.\cite{series:traceifaces,series:tieredobs}
\end{enumerate}
The deployer publishes a small \emph{public bundle} (receipt + digest-bound companion objects) and retains a \emph{private bundle} (raw logs, scripts, traces, and any privileged measurements).
Users do \emph{not} need the private bundle to understand or compare the claim; auditors can request it to replay the declared computations.\cite{series:auditorreplay}

\paragraph{Primary-path specificity without bloat.}
The public \texttt{primary\_surface\_manifest\_id} is intentionally a \emph{minimal guard object}: it commits to the separation class and the clauses that would invalidate the non-collusion assumption.\cite{series:primarysepmanifest}
Protocol- and endpoint-level details (e.g., the delegated-routing API version, endpoint set, service discovery/AutoConf mode, and any caching knobs) live in the private bundle as a config snapshot referenced by the manifest's evidence hooks; this keeps the public receipt diffable without forcing a full configuration dump into the paper.
For the deployed-surface taxonomy and a ``what-to-pin'' checklist (projection vs state vs separation), see Synthesis~30.\cite{series:deployedsurfaces}

\paragraph{Where the worked bundle lives.}
In the archive this example is represented by the receipt file \texttt{artifacts/example\_receipt.json} together with digest-bound companion objects.
The full file list is intentionally not reproduced here; it is discoverable from \texttt{support\_manifest.json} and \texttt{example\_artifact\_inventory.json} in the bundle.

A small family of adjuncts keeps the example auditable without bloating the note.
\texttt{example\_effective\_surface\_resolution.json} records how service-facing surfaces are assembled from the state declaration, contact surface, and ENF/projection.
\texttt{example\_release\_action\_matrix.json} and \texttt{example\_release\_obligation\_profile.json} separate ``what changed?'' from ``what must now be republished?''
\texttt{example\_successor\_status\_envelope.json} keeps the narrow public-status token distinct from the outward-facing lineage sentence that compresses it.
\texttt{example\_release\_spine.json} and \texttt{example\_release\_stage\_walkthrough.json} keep the stable release spine separate from one concrete release walk.

The downstream challenge layer is likewise split into small auditable owners rather than one omnibus note: routes, stop profiles, coverage, terminal witnesses, surface hooks, answer capsules, answer sheets, answer cards, carrier slots, audit trails, citation dockets, publication profiles, disclosure packets, escalation ladders, review menus, paper-facing public request contracts, requestable evidence classes, and request-fulfillment certificates each own exactly one late-stage task.
The maintained review-menu ledger, \texttt{example\_successor\_challenge\_answer\_review\_menus.json}, turns common referee follow-up requests (brief inline check, expanded inline check, exact-location check, full-audit check, request-packet handoff, or reveal-order check) into explicit smallest-sufficient handoff choices rather than leaving that response policy implicit in prose.
Its route/spine companions now also mirror the release-spine seam directly: repeated stage labels such as \emph{diff}, \emph{classify}, or \emph{summarize} recur there only as imported release-spine vocabulary, while the concrete stage order and stage-local stop rules remain attached to \texttt{example\_release\_spine.json}. They also mirror not only sentence-reuse and packet-reuse posture but the successor packet's family basis, support surface, manifest/inventory cuts, derived validation posture, and the visible-refresh branch's current closure status, closure scope, branch-local reopen triggers, and first reopen-owner handoffs, so a later terse note can distinguish ``the packet still refreshes in place'' from ``the refreshed visible branch is already a terminal public stop'', can see whether any shared validation companions remain outside that public basis, can still see what would reopen that stop, and can see which maintained owner regains authority first without reopening the full successor packet tail.
The companion evidence-class ledger, \texttt{example\_requestable\_evidence\_classes.json}, then gives stable names to the omitted-support families behind those request packets: bundle identity and validator companions, the public-status lineage/closure slice, and the compact-diff compare/replay slice. Its companion contract ledger now also exports a route-side public-anchor-family vocabulary map so the family keys reused by the evidence classes stay visibly owned by the paper-level contract rather than being silently flattened into later family or class mirrors. The new completion ledger, \texttt{example\_request\_fulfillment\_certificates.json}, says when those promised families have actually been fully delivered for the maintained worked bundle. The companion status, carryforward, notice, normal-form, selector, response-packet, packet-carryforward, delta-ledger, refresh-notice, refresh-notice-normal-form, refresh-notice-selection, refresh-response-menu, successor-facing refresh-response packet, and capstone closure ledgers now form an explicit eighteen-step source-only request ladder grouped into four adjacent owner buckets: the exact paper-level request-contract / evidence-class / fulfillment-certificate / status-envelope / carryforward-envelope / notice / canonical-notice-family / notice-normal-form / notice-selection chain, the exact response-menu / request-class-owner / response-packet / packet-scope / packet-family-basis / packet-support-surface / packet-manifest-cut / packet-inventory-cut / packet-carryforward-profile / packet-delta-ledger chain, the exact visible refresh-notice / canonical-refresh-family / refresh-normal-form / refresh-selection chain, and the exact successor-facing refresh-response-menu / successor-followup-owner / successor-facing packet / validation-posture / closure-verdict / closure-scope chain. That lets a source-only note say not only which maintained packet/menu cut stands behind a public sentence and what kind of omitted support that cut promises, but also where the first sufficient stop is for completion, visible notice, packetization, refresh, and closure without reopening the whole tail by prose.\cite{series:requestableevidenceclasses,series:requestfulfillmentcertificates,series:publicrequeststatusenvelopes,series:publicrequestcarryforward,series:publicrequestnotices,series:publicrequestnoticenormalforms,series:publicrequestnoticeselections,series:publicrequestresponsemenus,series:publicrequestresponsepackets,series:publicrequestresponsepacketcarryforward,series:publicrequestresponsepacketdeltas,series:publicrequestresponsepacketrefreshnotices,series:publicrequestresponsepacketrefreshnoticenormalforms,series:publicrequestresponsepacketrefreshnoticeselections,series:publicrequestresponsepacketrefreshresponsemenus,series:publicrequestresponsepacketrefreshresponsepackets,series:publicrequestresponsepacketrefreshresponsepacketclosureverdicts}

A new bridge adjunct, \texttt{example\_series\_spine.json}, now ties these layers together explicitly.
For each representative line item it names the mathematical root(s), receipt anchor, replay anchor, release-evolution anchor, and the answer/review anchor that later short papers should reuse.
For each shipped audience/question answer it then records only the stable cross-owner summary later terse papers actually need: the checked-sentence owners and inherited inline-provenance floor on the answer side, the release-obligation / closure / status chain in the middle, the paper-level request-contract / evidence-class / fulfillment / status / carryforward bridge that opens the source-only side, the bounded omitted-support packet and its packet-local basis / support-surface / manifest-cut / inventory-cut companions, the visible-refresh successor packet and closure verdict, and the paired terminal stop / widen / reopen tags that say whether a later paper may stop at the answer-side terminal form, widen to the adjacent request-side terminal form, or must reopen a named owner.
The exact packet ids, support-surface maps, manifest/inventory cuts, closure scopes, and reopen triggers are mirrored for audit but remain owned by the narrower notes and bundle ledgers rather than by the bridge itself.\cite{series:seriesspines,series:pairedterminalcitationdiscipline,series:pairedterminalcutoverrules}
Its companion \texttt{example\_question\_routes.json} then makes that bridge directly usable: for each common question family in \Cref{tab:workedexample-routingcuts} it records one concrete stem, the first sufficient stop, the full long-form walk when an auditor really wants every intermediate owner, and the shorter terminal forms / cutover tags when a later terse paper only needs the stable stop, adjacent widening capstone, or reopened owner.
That lets later notes cite one maintained theorem$\to$receipt$\to$replay$\to$release$\to$checked-sentence$\to$review bridge and then reopen only the first owner that the live question actually revives, instead of silently restitching the whole tail.\cite{series:seriesspines,series:pairedterminalcitationdiscipline,series:pairedterminalcutoverrules}

\begin{table}[t]
\centering
\small
\begin{tabularx}{\linewidth}{@{}Y Y@{}}
\toprule
Public object (digest-bound) & What it fixes (why a user can diff it)\\
\midrule
\texttt{receipt\_id} & The line items: declared interfaces, budgets, and replay hooks (Synthesis~12).\\
\texttt{tw\_id} & Threat window declaration: secret, repetition unit/window, adversary model (Synthesis~4).\\
\texttt{enf\_registry\_id} & Exposure normal form labels: which transcript projection is in scope (Synthesis~3).\\
\texttt{state\_decl\_id} & The small public state surface that conditions the claim (state series + Synthesis~18).\\
\texttt{primary\_surface\_manifest\_id} & Guard object for non-collusion/separation assumptions (Synthesis~22).\\
\texttt{replay\_plans\_id} & Auditor replay recipe: how to recompute each line item from private artifacts (Synthesis~20).\\
\bottomrule
\end{tabularx}
\caption{Minimal public identifiers for the worked example. The point is not the particular files but the \emph{digest discipline}: a name like \texttt{enf.fallback.cct.v1} is not trusted unless its digest matches.}
\label{tab:min-public-ids}
\end{table}

\section{Receipt line items}
\label{sec:lineitems}
This note does not re-specify the receipt schema; it \emph{instantiates} it.
The canonical tuple schema lives in Synthesis~12, and deployed control-plane surfaces referenced by any line item are owned by Synthesis~30.\cite{series:receiptitems,series:deployedsurfaces}
Here we keep only the slice needed to make the scenario mechanically replayable.

\paragraph{Minimal slice used here.}
\begin{enumerate}[leftmargin=*]
\item \textbf{Interface IDs:} \texttt{tw\_id} (threat window template) and \texttt{exposure\_nf\_id} (declared observable projection / ENF)---owned by Synthesis~4 and Synthesis~3.\cite{series:threatwindows,series:traceifaces}
\item \textbf{State binding:} \texttt{state\_decl\_id} pins the state surfaces this bundle admits (e.g., the public refresh timer and bucket-level staleness) and therefore what ``replay'' means.\cite{series:planstateequiv}
\item \textbf{Primary separation:} \texttt{primary\_path\_declaration} publishes the separation class and its failure guard $(0,\rho)$.\cite{series:primarysepmanifest,series:probsep}
\item \textbf{Fallback exposure:} \texttt{fallback\_contact\_surface} and \texttt{fallback\_tiered\_vector} publish attenuated Tier~2 summaries under an explicit activation cap.\cite{series:obsatten,series:fallbacktax,series:tieredobs}
\item \textbf{Routing-signature compression route:} Certified~B owns the exact raw-$M$ to $M_{\mathrm{eff}}$ compression card for the declared committee-contact projection and weekday/weekend context support. In this worked cut, \texttt{routing\_signature\_manifest} is only an owner-map / series-spine boundary row with \texttt{materialization\_status=owner\_map\_spine\_only\_current\_cut}; it is not a receipt line item, UVI/verifier packet, replay-plan row, or support-bundle route.\cite{series:certifiedB}
\item \textbf{Effective surface assembly:} the worked bundle ships a tiny owner-side adjunct showing the three-way split \texttt{state\_decl\_id} vs. \texttt{contact\_surface\_id} vs. \texttt{exposure\_nf\_id}; this is where replay-time inheritance and precedence are made explicit without bloating the receipt schema.\cite{series:receiptitems,series:deployedsurfaces}
\item \textbf{Repeated-use profiling:} \texttt{profiling\_equalization} publishes a profiling/equalization summary that can be compared across releases.\cite{series:profiling}
\end{enumerate}

\begin{table}[ht]
\centering
\small
\begin{tabularx}{\linewidth}{@{}Y Y Y Y@{}}
\toprule
Audit question & First artifact to inspect & Stop here when & Widen only if \\
\midrule
What public claim is actually being made? & \texttt{example\_receipt.json} & You only need the released line items, public identifiers, and digest-bound claim surface. & The question becomes one of theorem ownership, replay, release drift, or follow-up packaging rather than claim identification. \\
Which owner note or mathematics backbone justifies that line item? & \texttt{example\_line\_item\_owner\_map.json} (+ Synthesis~24 when a published root is involved) & You only need the receipt-facing owner and, when relevant, the canonical mathematical root for that line item. & The question becomes where the claim lands publicly, how it is replayed, or how it evolves across releases. \\
How would an auditor recompute the number? & \texttt{example\_replay\_plans.json} & You only need the replay plan ids, required private inputs, and plan-spec hooks for recomputation. & The question becomes which maintained release stage, archive cut, or referee handoff moves after a drift. \\
What moves if the claim changes across releases? & \texttt{example\_release\_spine.json} together with \texttt{example\_compare\_report.json} & You only need the change classification and the maintained stage family that now owns the next answer. & The question becomes the exact public-on-request archive cut or a later concrete follow-up class. \\
What exact archive slice should travel when omitted support is requested? & \texttt{example\_successor\_challenge\_answer\_disclosure\_packets.json} & You only need the exact requestable archive slice and its manifest/inventory/validator companions for that one public answer. & The question widens from ``what slice travels?'' to ``which concrete follow-up class is smallest sufficient after that slice?'' or leaves the answer side for the source-only request tail. \\
What is the smallest sufficient maintained follow-up handoff after that? & \texttt{example\_successor\_challenge\_answer\_review\_menus.json} & You only need the smallest sufficient handoff for one concrete referee request class after the paper-facing archive cut. & The question stops being answer-side follow-up packaging and becomes a source-only request/refresh-packet owner question. \\
\bottomrule
\end{tabularx}
\caption{The worked example becomes referee-usable when the reader follows the first sufficient audit stop and widens only when the question actually changes.}
\label{tab:workedexample-auditwalk}
\end{table}
\paragraph{Stable ladder extracted from the audit walk.}
For each representative line item, the maintained left-to-right path is now explicit: theorem or imported mathematical root $\to$ receipt-facing owner $\to$ replay plan $\to$ release-stage owner $\to$ answer-side disclosure packet (when the task is ``what exactly do we hand over on request?'') $\to$ review handoff.
The practical rule is now exported mechanically as well as rhetorically: stop at the first sufficient owner for the question actually being asked, and widen only when the question itself changes.
Synthesis~24 owns the theorem-root side of that ladder, Synthesis~32 owns the release-local stage cuts, Synthesis~52 owns the public-on-request disclosure cut, Synthesis~55 owns the cross-series bridge, Synthesis~76 owns the wider theorem-to-publication stop ladder above those local owners, and this worked note owns only the concrete instantiation of the maintained worked stop beneath that wider ladder.\cite{series:mathcrosswalk,series:releasespine,series:challengeanswerdisclosurepackets,series:seriesspines,series:theoremtopublicationspine}

\paragraph{A worked publication-boundary card.}
One recurring ambiguity in terse papers is no longer mathematical or replay-local at all; it is whether the live question is still one concrete worked cut below publication, has crossed into the exact on-request/source-only tail, or has widened back above this worked route into the broader theorem-to-publication stop ladder. The worked bundle is now explicit enough to answer that question with exact local artifacts instead of wrapper prose.\cite{series:challengeanswerdisclosurepackets,series:challengeanswerreviewmenus,series:publicrequestresponsepacketrefreshresponsepacketclosureverdicts,series:pairedterminalcutoverrules,series:theoremtopublicationspine}

\begin{table}[ht]
\centering
\small
\begin{tabularx}{\linewidth}{@{}Y Y Y@{}}
\toprule
Live question at the publication boundary & First honest worked stop & Widen only if \\
\midrule
What exact archive slice sits behind the shipped checked answer? & \texttt{example\_successor\_challenge\_answer\_disclosure\_packets.json} & The question becomes which concrete referee follow-up class is smallest sufficient after that slice, or has genuinely crossed into the source-only request branch. \\
Which smallest sufficient answer-side follow-up handoff now answers the referee? & \texttt{example\_successor\_challenge\_answer\_review\_menus.json} & The live issue stops being answer-side follow-up packaging and becomes either one source-only request/refresh owner question or the paired stop/widen/reopen judgment itself. \\
Has the live issue genuinely crossed the archive cut into the first bounded source-only packet? & \texttt{worked\_example\_receipt\_interlock\_successor\_public\_request\_response\_menu\_\_request\_scope\_check} & The bounded packet is no longer enough and the live issue has become the exact successor-facing visible-refresh packet or its closure verdict. \\
Where does the adjacent request-side visible-refresh branch currently stop? & \texttt{worked\_example\_receipt\_interlock\_successor\_public\_request\_response\_menu\_\_request\_scope\_check\_\_delta\_\_refresh\_notice\_\_selection\_\_response\_menu\_\_packet\_membership\_check\_\_closure} & A named request-side reopen trigger has fired, or the live issue is no longer one request-side terminal question but the wider theorem/worked/release/notary/claim/release stop selection. \\
Which wider theorem/worked/release/notary/claim/release owner should carry the claim at all? & Synthesis~76 & The live issue narrows again to one concrete worked cut, in which case reopen the first sufficient worked stop above rather than keeping the wider ladder open. \\
\bottomrule
\end{tabularx}
\caption{A worked publication-boundary card. This is the smallest concrete map from one shipped checked answer to its exact answer-side packet, review-side handoff, first bounded source-only packet, request-side terminal closure, and the wider theorem-to-publication ladder above them.}
\label{tab:worked-publication-boundary-card}
\end{table}

\paragraph{One tiny publication-boundary drill.}
Suppose a later terse paper starts with the shipped checked answer and the public sentence ``supporting artifacts are available on request.'' The first honest worked question is usually not ``open the whole tail'' but ``what exact archive slice travels?'' so the paper stops first at the disclosure packet. If the referee instead asks which concrete follow-up class is smallest sufficient after that slice, the paper stops next at the review menu. If the live issue then genuinely leaves the answer side, the paper may widen to the first bounded source-only packet and, only if the visible-refresh closure itself is now the live issue, to the exact request-side terminal closure verdict. If the dispute is instead whether this claim should stop at one worked cut, one replay/notary owner, one claim object, or one lineaged release at all, the honest widening target is Synthesis~76 rather than more worked-tail prose.\cite{series:challengeanswerdisclosurepackets,series:challengeanswerreviewmenus,series:publicrequestresponsepacketrefreshresponsepacketclosureverdicts,series:pairedterminalcutoverrules,series:theoremtopublicationspine}
\paragraph{Concrete seam versus abstract tail owners.}
This worked publication-boundary card owns only one concrete maintained seam. Later terse papers should stop here when they need the exact audited route from the shipped checked answer into its adjacent request tail for this worked branch. If they instead need the abstract claim that the answer-side terminal stop and the adjacent request-side capstone may travel together as one paired sentence, they should reopen Synthesis~74. If they need the abstract stop/widen/reopen rule rather than this one maintained branch, they should reopen Synthesis~75. If they need wider stop selection above the publication boundary, they should reopen Synthesis~76 rather than stretching this worked card into a route owner.\cite{series:pairedterminalcitationdiscipline,series:pairedterminalcutoverrules,series:theoremtopublicationspine}

\paragraph{One congestion owner-map boundary drill.}
The current worked-example cut exposes the congestion-family route through \nolinkurl{example\_line\_item\_owner\_map.json} and the generated series spine rather than through receipt/support/verifier packet rows. The owner map includes \nolinkurl{congestion\_epoch\_contract} and \nolinkurl{pscq\_mechanism\_packet}, and it keeps W-Congestion-EQ adjacent as the replay owner for the same family route. However, \nolinkurl{example\_receipt.json} does not contain \nolinkurl{congestion\_epoch\_contract} or \nolinkurl{pscq\_mechanism\_packet}; \nolinkurl{example\_support\_bundle\_map.json} does not contain \nolinkurl{bundle://worked-example/congestion-eq} or \nolinkurl{bundle://worked-example/pscq}; and \nolinkurl{example\_verifier\_report.json} does not contain \nolinkurl{congestion\_replay\_packet}. The honest drill is therefore a boundary drill, not a packet drill: stop at Congestion~1 for the local waiting-time witness/load-envelope/accountant card, add PSC-Q for the local calendar/substitution mechanism card, and widen to W-Congestion-EQ for the local replayable checker card; use the worked-example owner map and Synthesis~55 only to confirm the route and owner split.\cite{series:congestion1,series:pscq,series:wcongestion,series:seriesspines}

\paragraph{One routing-signature owner-map boundary drill.}
The current worked-example cut exposes the certified-routing compression route through \nolinkurl{example\_line\_item\_owner\_map.json} and the generated series spine rather than through receipt/support/verifier packet rows. The owner map includes \nolinkurl{routing\_signature\_manifest} and names Certified~B as the first public owner for the compression witness, with Certified~A and Certified~C adjacent as theorem-root and remanifest/recertify neighbors. However, \nolinkurl{example\_receipt.json} does not contain a \nolinkurl{routing\_signature\_manifest} line item; \nolinkurl{example\_replay\_plans.json} does not contain \nolinkurl{eval-routing-signature-v1}; and \nolinkurl{example\_support\_bundle\_map.json} does not contain \nolinkurl{bundle://worked-example/routing-signature}. The honest drill is therefore a boundary drill, not a packet drill: stop at Certified~B for the local routing-signature quotient map and effective-multiplicity witness, widen to Certified~C only for support drift or recertification, and use the worked-example owner map plus Synthesis~55 only to confirm that owner split.\cite{series:certifiedmenus,series:certifiedB,series:recert,series:seriesspines}

\paragraph{One recertification source-object boundary drill.}
The maintained worked bundle exposes the concrete Certified~C change-control drill through source objects, not through a separate recertification support-bundle packet.
\nolinkurl{example\_change\_control.json} fixes the same-claim guard as \nolinkurl{change\_control\_id=chgctl.receiptchain.v1}, stable \nolinkurl{tw\_id=sha256:c1ff81fc84049a2160a21487f61a4565e03c5263f1437649c4299d1354b35c18}, stable \nolinkurl{state\_decl\_id=sha256:ddfbb55199a6ecdb225cbafc60a74c234ae3334fb40528fde0a2ec2ce6de3f20} for \nolinkurl{stateadj.publicbucket.v1}, unchanged required exposures, inherited margin \nolinkurl{gamma=0.020}, and primary separation class \nolinkurl{sep.nr1}.
\nolinkurl{example\_compare\_report.json} then records the concrete successor diff: the guard fields stay fixed, \nolinkurl{fallback\_contact\_surface} and \nolinkurl{fallback\_tiered\_vector} change, the maintained classification is \nolinkurl{replay\_changed\_surfaces}, and the required replay scope is exactly \nolinkurl{eval-fallback-cct-v1} plus \nolinkurl{eval-fallback-vector-v1}.
\nolinkurl{example\_drift\_cases.json} together with \nolinkurl{example\_release\_action\_matrix.json} makes the cutover ladder exact: \nolinkurl{refresh-only-A} stays wrapper-only, \nolinkurl{replay-required-B} lands at rule \nolinkurl{contact-or-budget-replay}, and \nolinkurl{full-recertification-C} lands at \nolinkurl{interface-or-state-recertify}.
The boundary is equally important: \nolinkurl{example\_support\_bundle\_map.json} does \emph{not} resolve \nolinkurl{bundle://worked-example/recertification}, \nolinkurl{example\_replay\_plans.json} does not include \nolinkurl{eval-recertification-v1}, and \nolinkurl{example\_verifier\_report.json} emits no separate recertification packet.
So the worked stop is concrete as a source-object boundary: a later terse paper may stop at Certified~C when the live issue is whether the successor stayed under the same certified interface and which side of the replay/recertify cut it landed on; widen to Operational~B only when the live issue becomes live deployment conformance, and widen to Eval~1 or Release~A only when the live issue becomes notarized comparison or signed public lineage above that boundary.\cite{series:recert,series:operationalB,series:eval1,series:releaseA}

\paragraph{One exact profiling-contract drill.}
The maintained worked bundle also carries one exact anonymous-DHT profiling stop rather than only a generic label.
In \nolinkurl{example\_receipt.json}, the line item \nolinkurl{profiling\_equalization} fixes \nolinkurl{exposure\_nf\_id=sha256:be229915915d36aed50024b2ed38da0d781be75bd044a5b9d70287808b356798}, \nolinkurl{tw\_id=sha256:c1ff81fc84049a2160a21487f61a4565e03c5263f1437649c4299d1354b35c18}, \nolinkurl{budget\_value=(eta\_bits=0.821760, delta=0.007464)}, \nolinkurl{projection=retry\_bucket\_trace}, \nolinkurl{bucket\_width=250ms}, \nolinkurl{evidence\_id=sha256:7ded7411c77f80f2b3a674baff4e137807e3c1f769b2399ff87588b62b9c3cad}, and replay hook \nolinkurl{eval-profiling-eq-v1}/\nolinkurl{sha256:db245ce4f3c200322eea17a42a2c64cf9dc28bdeb18fa24e5d3643b51eb1e6f6}.
\nolinkurl{example\_line\_item\_owner\_map.json} then names Anonymity~B as the first public owner for that packet and points to \nolinkurl{example\_profiling\_evidence.json} plus \nolinkurl{example\_replay\_plans.json} as the public objects attached to it.
\nolinkurl{example\_support\_bundle\_map.json} resolves \nolinkurl{bundle://worked-example/profiling-eq} to \nolinkurl{example\_receipt.json}, \nolinkurl{example\_profiling\_evidence.json}, and \nolinkurl{support://profiling/log-slice.example}.
So the worked stop is now concrete: a later terse paper may stop at the profiling card when the live issue is the declared retry-bucket equalization witness itself; widen to Synthesis~20 only when the live issue becomes the exact replay checklist or notarized recomputation, and widen to ABOM/OINL or Release~A only when the live issue becomes the signed claim object or lineage wrapper above that card.\cite{series:profiling,series:replayrecipe,series:abomoinl,series:releaseproperty}


\paragraph{One Boss Fight dial-sheet owner-map boundary drill.}
The current worked-example cut exposes the Boss Fight dial-sheet route through \nolinkurl{example\_line\_item\_owner\_map.json} and the generated series spine rather than through a new receipt/support packet. The owner map includes \texttt{bossfight\_dialsheet\_packet} and keeps Boss Fight~A as the repeated-use accountant root, Boss Fight~B as the knob-surface owner, the Boss Fight~B addendum as optional representative evidence, and Boss Fight~C as the replay companion. However, \texttt{example\_receipt.json} does not contain \texttt{bossfight\_dialsheet\_packet}; \texttt{example\_replay\_plans.json} does not contain \texttt{eval-bossfight-dialsheet-v1}; and \texttt{example\_support\_bundle\_map.json} does not contain \texttt{bundle://worked-example/bossfight-dialsheet}. The honest drill is therefore a boundary drill, not a packet drill: stop at Boss Fight~B for the local dial-sheet card, widen to the addendum only for representative support tables, and widen to Boss Fight~C only for the replayable checker contract.\cite{series:bossfightB,series:bossfightBaddendum,series:bossfightC,series:evalnotary,series:releaseproperty,series:workedexample,series:seriesspines}

The transcendental endpoint derivation bundle checks the maintained log2 endpoint upper bounds by exact rational series and integer witnesses, while \nolinkurl{profiling_statistical_derivation_bundle} now proves the base profiling eta/delta CP+union-bound envelope by exact binomial-tail/log-ratio checks under \nolinkurl{fixed_sample_stop_rule=worked-example-profiling-fixed-sample-lock-v1}, the fixed-sample stop/refresh lock.  The verifier treats that row as \nolinkurl{not_time_uniform_fixed_sample_only}: one frozen count-table digest, one allowed look, no peeking or repeated refresh reuse, and a successor evidence ID or true time-uniform certificate before the profiling evidence can be refreshed.

\paragraph{One operational owner-map boundary drill.}
The current worked-example cut exposes the retry and runtime-conformance operational route through \nolinkurl{example\_line\_item\_owner\_map.json} and \nolinkurl{example\_series\_spine.json}, not through emitted packet rows. The owner map includes \nolinkurl{retry\_schedule\_packet} and \nolinkurl{runtime\_conformance\_packet}; the series spine carries owner-map / series-spine ladders for those names so later terse papers can see the intended Operational~A -> Operational~B split. However, \nolinkurl{example\_receipt.json} does not contain either line item, \nolinkurl{example\_replay\_plans.json} does not contain \nolinkurl{eval-retry-schedule-v1} or \nolinkurl{eval-runtime-conformance-v1}, and \nolinkurl{example\_support\_bundle\_map.json} does not contain \nolinkurl{bundle://worked-example/retry-schedule} or \nolinkurl{bundle://worked-example/runtime-conformance}. The honest drill is therefore a boundary drill, not a packet drill: stop at Operational~A for the local schedule-only retry theorem/card, widen to Operational~B for local pinned-plan/PTL/epoch-conformance evidence, and use the worked-example owner map plus Synthesis~55 only to confirm the route split.\cite{series:profiling,series:schedretries,series:operationalB,series:bossfightA,series:advantagecontracts,series:releaseproperty,series:workedexample,series:seriesspines}

\paragraph{One Boss Fight verifier owner-map boundary drill.}
The current worked-example cut exposes the Boss Fight verifier route through \nolinkurl{example\_line\_item\_owner\_map.json} and the generated series spine rather than through a new \texttt{example\_verifier\_report.json} packet. The owner map includes \texttt{bossfight\_verifier\_bundle} and keeps Boss Fight~A, Boss Fight~B, and Boss Fight~C as the linear accountant / dial-sheet / checker owners, with Eval~1 and Release~A downstream. However, \texttt{example\_verifier\_report.json} does not contain \texttt{bossfight\_verifier\_packet}; \texttt{example\_replay\_plans.json} does not contain \texttt{eval-bossfight-verifier-v1}; and \texttt{example\_support\_bundle\_map.json} does not contain \texttt{bundle://worked-example/bossfight-verifier}. The honest drill is therefore a boundary drill, not a packet drill: stop at Boss Fight~C for the local replayable checker object, widen to Eval~1 only for attempted-evaluation notarization, and widen to Release~A only for closure packaging above that local checker.\cite{series:bossfightC,series:bossfightBaddendum,series:evalnotary,series:releaseproperty,series:workedexample,series:seriesspines}

\section{What changes across a release (micro-diffs)}

This worked bundle is meant to make \emph{semantic drift mechanically visible}.
Three common drift patterns illustrate what should (and should not) change in a receipt:

\begin{itemize}[leftmargin=*]
\item \textbf{Knob drift (same interface).} If the transcript projection is unchanged but a retry/timeout/width knob changes, the line item keeps the same \texttt{exposure\_nf\_id} but changes \texttt{knobs} and (usually) \texttt{budget\_value}. Tag \texttt{lineage=same-enf} and require replay/evidence for the new number.\cite{series:receiptitems,series:planstateequiv}
\item \textbf{State-surface drift (same mechanism, different reader set).} If a deployment silently changes which delegated-routing endpoints are contacted (AutoConf defaults, cache layers, gateway exposure), the \texttt{state\_decl\_id} must change even if the budget arithmetic is identical. Otherwise the claim is underspecified.\cite{series:deployedsurfaces,series:planstateequiv}
\item \textbf{Interface drift (new projection).} If the observable transcript grows (e.g., a new \texttt{/routing/v1} endpoint is used, response streaming changes message timing, or ``empty result'' behavior changes from a status-coded 404 surface to a 200-empty compatibility surface), treat it as a new \texttt{exposure\_nf\_id} (or an explicit coarsening if the projection is strictly reduced). Do not pretend the previous certificate applies.\cite{series:traceifaces,series:receiptitems,series:deployedsurfaces}
\end{itemize}

The equivalence/coarsening taxonomy for plan/state ids in Synthesis~18 is the intended shared language for these cases; this section is only the ``worked example'' view.\cite{series:planstateequiv}

\paragraph{User micro-diff algorithm (receipt-level).}
A user does not need to interpret every field; a deterministic diff is enough:
(1) fetch both receipts and verify their advertised digests (\texttt{receipt\_id});
(2) compare the small set of top-level interface/state ids;
(3) compare line items by \texttt{name}.
\begin{verbatim}
Top-level: tw_id, enf_registry_id, state_decl_id,
           primary_surface_manifest_id, replay_plans_id
Per line item (by name): exposure_nf_id, contact_surface_id, knobs,
                         budget_value, evidence_id, replay_hook, lineage
\end{verbatim}
Interpretation: changing \texttt{exposure\_nf\_id} is interface drift; changing \texttt{state\_decl\_id} is control-plane/reader-set drift; changing \texttt{knobs} or \texttt{budget\_value} with the same \texttt{exposure\_nf\_id} is knob drift (same interface, new certificate).

\paragraph{From diff to action.}
The shipped bundle now includes \texttt{artifacts/example\_release\_action\_matrix.json}, a digest-bound field-family-to-action matrix that instantiates Synthesis~18 for this one receipt chain.
Its job is triage, not new theory: with stable \texttt{tw\_id}/\texttt{state\_decl\_id}/required \texttt{exposure\_nf\_id}s but changed \texttt{contact\_surface\_id}, \texttt{knobs}, \texttt{budget\_value}, or replay-plan digests, the default action is ``replay the changed surfaces and publish a fresh notarized comparison''; if \texttt{tw\_id}, \texttt{state\_decl\_id}, a required \texttt{exposure\_nf\_id}, or the primary separation guard changes, continuity claims escalate to full recertification; pure log-anchor or watch-text updates remain refresh-only. The sibling \texttt{example\_release\_obligation\_profile.json} then answers the separate publication question by listing the minimal refreshed outputs for each action class (for Case~B: fresh release binding, compare report, replay verdict, successor continuity verdict, and successor lineage notice). This keeps the worked note tight while giving users one checkable answer to both ``what do I do now?'' and ``what should I expect to be republished?''.

\paragraph{Canonicalization notes.}
Digest IDs should be computed over canonical bytes, not ad hoc pretty-printed JSON.
For JSON objects we recommend the JSON Canonicalization Scheme (JCS, RFC~8785), i.e., define \texttt{sha256(JCS(obj))}.\cite{rfc8785}
For human micro-diffs between two receipts-as-JSON, also canonicalize \emph{semantic sets} (sort and deduplicate arrays that represent sets, such as protocol filter lists or router allowlists) and normalize duration strings (parse \texttt{30s}/\texttt{0.5m} into a single unit) so reordering/formatting noise does not masquerade as drift.

\paragraph{Canned drift cases (bundle-provided).}
The shipped bundle includes \texttt{artifacts/example\_drift\_cases.json}, a three-row catalog of ``what changed'' successors for the current base release label \texttt{worked-example-draft-189}. It is intended as a sanity check for the taxonomy above and for the shipped release-action matrix (refresh-only vs replay-required vs recertification).
\begin{table}[t]
\centering
\begin{tabularx}{\linewidth}{@{}Y Y Y@{}}
\toprule
Case & Class & What changed / required action \\
\midrule
\texttt{refresh-only-A} & \texttt{refresh\_only} & Log checkpoint / user-watch text only; rebind receipt lineage, no budget re-evaluation. \\
\texttt{replay-required-B} & \texttt{replay\_changed\_surfaces} & Fallback observation upper bound tightened ($p_{\mathrm{obs}}:0.02\to0.015$); rerun declared evaluators and publish a fresh notarized comparison. \\
\texttt{full-recertification-C} & \texttt{full\_recertification} & Threat window and state declaration (cache trigger) changed; treat as certified-interface drift and require recertification gating. \\
\bottomrule
\end{tabularx}
\caption{Three canned drift cases shipped in the worked-example bundle.}\label{tab:workedexample-driftcases}
\end{table}

\paragraph{Example output (from the shipped compare report).}
The bundle also includes \texttt{artifacts/example\_compare\_report.json}; its \texttt{user\_fast\_diff} contains a ``first alarm'' message and the minimal numeric fields that casual users usually care about. For the canned successor in that report, the fallback budget improves while \texttt{exposure\_nf\_id} stays fixed:
\begin{verbatim}
base:     budget_value=0.05635, total_summary_bits=0.26590
successor budget_value=0.04240, total_summary_bits=0.26518
(same exposure_nf_id; receipt digest changed)
\end{verbatim}

\paragraph{Derived comparison verdict (bundle-provided).}
The bundle now also ships \texttt{artifacts/example\_successor\_continuity\_verdict.json}, a tiny base-to-successor capsule that records the matched action-matrix rule, the replay scope (here \texttt{fallback\_contact\_surface} and \texttt{fallback\_tiered\_vector}), and the continuity decision \texttt{same\_certified\_interface\_replayed\_surfaces}. It is intentionally not a new receipt primitive: it is the smallest checkable answer to ``does continuity still hold for this specific successor, and what had to be replayed to say so?''

\paragraph{Derived closure ledger and publication-closure verdict (bundle-provided).}
The same canned comparison now also ships \texttt{artifacts/example\_release\_closure\_ledger.json}, a tiny role-by-role bridge between the obligation profile and the closure verdict. It does not reclassify the diff or re-run the replay. Its narrower job is to list, for the winning obligation class, each required publication role, the allowed concrete artifact candidates named by the profile, the artifact(s) actually bound in this shipped package, and whether that row is satisfied. The companion \texttt{artifacts/example\_publication\_closure\_verdict.json} then compresses that ledger to the outward answer: did this concrete successor package actually publish the outputs required by the winning obligation class? In the worked Case~B path, the ledger and verdict bind \texttt{fresh\_release\_binding}, \texttt{fresh\_compare\_report}, \texttt{fresh\_replay\_verdict}, \texttt{fresh\_successor\_continuity\_verdict}, and \texttt{fresh\_successor\_lineage\_notice} to the concrete shipped artifacts. That keeps publication policy, row-level closure checking, and the final completeness token separate enough that later papers can cite only the layer they actually need.\cite{series:releaseclosure}

\paragraph{Derived successor-lineage notice (bundle-provided).}
The bundle now also ships \texttt{artifacts/example\_successor\_lineage\_notice.json}, a tiny outward-facing status capsule for the same canned comparison. It does not replace the compare report, continuity verdict, or closure verdict; it simply collates their public answer into one citeable object: this successor stays on the same certified interface, the replay scope is the fallback contact/vector slice, and the required successor publication outputs are complete. That gives later papers and release notes one compact status object to cite, while readers who want details can follow its pointers back to the narrower maintenance adjuncts.

\paragraph{Derived successor-derivation graph (bundle-provided).}
The bundle now also ships \texttt{artifacts/example\_successor\_derivation\_graph.json}, a tiny maintenance DAG for the same canned comparison. It does not add a new verdict. Its job is narrower: make the dependency order explicit so later notes can cite one graph rather than re-explaining that the compare report and replay verdict feed the continuity verdict, the obligation profile and compare report feed the closure ledger, the closure ledger feeds the closure verdict, the public lineage notice is only a wrapper on top of those narrower objects, and the claim-support map is last because it merely binds outward-facing sentences back to that already-derived stack. That keeps maintenance prose short while still letting a referee verify the order in which the successor conclusion was assembled.

\paragraph{Derived successor-claim support map (bundle-provided).}
The bundle now also ships \texttt{artifacts/example\_successor\_claim\_support\_map.json}, \texttt{artifacts/example\_successor\_citation\_map.json}, \texttt{artifacts/example\_successor\_challenge\_routes.json}, \texttt{artifacts/example\_successor\_quote\_map.json}, and \texttt{artifacts/example\_successor\_clause\_pack.json}, a five-object downstream suffix for the same canned comparison.
Read left-to-right, they answer progressively narrower questions: which clauses are supported, what is the smallest sufficient thing to cite, where a release note/referee/auditor should start and how that route escalates if challenged, which exact field should be quoted, and which short reusable sentence is already checked.
The maintained bundle now also ships \texttt{artifacts/example\_successor\_normal\_form.json}, a tiny wrapper on just that suffix.
Its job is not to invent another verdict.
It simply freezes the summarize $\to$ support $\to$ cite $\to$ quote $\to$ clause path in one auditable object so later papers can cite the downstream normalization layer without re-listing the five adjuncts.
Instead of re-listing the whole late-maintenance tail every time, the worked bundle now treats it as seven owner families, with the source-only tail formalized as an explicit eighteen-step request ladder grouped into four adjacent owner buckets. Synthesis~31--33 own the release-local suffix and closure bridge. Synthesis~36--44 own challenge routing, stop profiles, coverage certificates, surface hooks, and terminal-witness discipline. Synthesis~45--54 own the assembled answer bundles and their handoff surfaces (capsules, sheets, cards, carrier slots, audit trails, citation/provenance profiles, disclosure packets, escalation ladders, and review menus). Within that answer-side bucket the worked bundle now also makes one additional seam explicit: repeated six-class review-request labels may recur across later routes and series-spine mirrors only as imported review-menu vocabulary, while the concrete handoff choice for each label remains owned by the audience/question review menu. Synthesis~56--63 own within-release source-only request promises (contracts, evidence-class names, completion criteria, status/carryforward tokens, outward notices, notice families, and family selection). Within that bucket the worked bundle now makes five seams explicit. First, repeated public-anchor-family keys may recur across evidence-class ledgers, fulfillment certificates, routes, and series-spine mirrors only as imported contract vocabulary, while the owner of each family's public-path / review-menu / disclosure-packet basis remains the paper-level public request contract. Second, repeated evidence-class keys may recur across later contracts, fulfillment certificates, routes, and series-spine mirrors only as imported class vocabulary, while the owner of each class's path-family meaning remains the evidence-class ledger itself. Third, the bundle-identity/validator-companion evidence class is shared across both public anchor families, but completion still lives in the two family-scoped fulfillment certificates rather than in a freestanding cross-family verdict. Fourth, repeated family-completion verdict names may recur across later status envelopes, visible notices, routes, and series-spine mirrors only as imported family-scoped certificate vocabulary, while the truth of those repeated verdict names still belongs only to the concrete fulfillment certificates that emit them. Fifth, the paper-level status and carryforward tokens are shared basis vocabulary that may recur across many visible notices, canonical notice families, selectors, routes, and series-spine mirrors, but those repeated token names still derive their truth only from the concrete status and carryforward envelopes rather than from the visible wrappers that import them. Synthesis~64--67 own notice-keyed response menus and bounded response packets together with their carryforward and delta rules. Within that bucket the worked note now makes two seams explicit. First, the refreshable-field-family labels in packet carryforward profiles are shared packet-maintenance vocabulary that may recur across many reusable packets, but packet reuse verdicts remain packet-scoped and are merely mirrored by later owners. Second, the changed-field-family labels in the packet delta ledgers are shared packet-local vocabulary that may recur across many refreshed packets, but the exact refresh instruction rows remain attached to concrete owner path/key pairs and they do not by themselves choose any visible wrapper family. The worked route/spine mirrors now also say directly that repeated canonical refresh-family labels recur there only as imported packet-local delta vocabulary rather than as new route-owned wrapper policy. Synthesis~68--70 own the visible successor-refresh layer for reused packets, where emitted sentences and canonical family choice remain notice-scoped even when they compress that shared delta vocabulary. Synthesis~71--73 own the successor-facing refresh-response menu, exact successor-facing packet, and capstone closure verdict for that reissued visible branch. The worked note now makes one final seam explicit there as well: a successor-facing packet may in principle reuse shared validation companions, but that support posture remains packet-local support surface rather than selector logic, and closure is certified only for the exact public packet basis. The companion route bridge now therefore keeps the exact answer-carrier-slot key and exact answer-audit-trail key visible beside the exact citation-docket key, exact disclosure-packet key, the exact source-only request-contract key, exact family-scoped fulfillment-certificate key list, exact source-only status-envelope key, exact source-only carryforward-envelope key, exact source-only notice key, exact source-only canonical-notice-family / notice-normal-form key, exact predecessor packet-carryforward-profile key, exact packet-delta key, exact successor refresh-notice key, exact successor canonical-refresh-family / refresh-normal-form key, exact refresh-notice-selection key, exact refresh-response-menu key, exact successor-facing packet key, and exact closure-verdict key visible beside the mirrored reuse / validation / closure tags rather than compressing that owner chain into posture labels alone. In the current worked cut every successor-facing packet is public-only, so the mirrored validation-posture map stays trivial while the mirrored closure-scope map stays equal to the public packet basis itself. This worked note instantiates those families once so the audit path is concrete without turning the example into a second index note.\cite{series:downstreamnormalforms,series:releasespine,series:releaseclosure,series:challengeroutes,series:challengestops,series:challengebranches,series:challengecoverage,series:challengesurfacehooks,series:challengeterminalwitnesses,series:challengecoveragewitnesspacks,series:challengecoveragewitnessslices,series:challengecoveragewitnesscores,series:challengeanswercapsules,series:challengeanswersheets,series:challengeanswercards,series:challengeanswercarrierslots,series:challengeansweraudittrails,series:challengeanswercitationdockets,series:challengeanswerpublicationprofiles,series:challengeanswerdisclosurepackets,series:challengeanswerescalationladders,series:challengeanswerreviewmenus,series:publicrequestcontracts,series:requestableevidenceclasses,series:requestfulfillmentcertificates,series:publicrequeststatusenvelopes,series:publicrequestcarryforward,series:publicrequestnotices,series:publicrequestnoticenormalforms,series:publicrequestnoticeselections,series:publicrequestresponsemenus,series:publicrequestresponsepackets,series:publicrequestresponsepacketcarryforward,series:publicrequestresponsepacketdeltas,series:publicrequestresponsepacketrefreshnotices,series:publicrequestresponsepacketrefreshnoticenormalforms,series:publicrequestresponsepacketrefreshnoticeselections,series:publicrequestresponsepacketrefreshresponsemenus}

\paragraph{A concrete coverage-to-capsule spine.}
One place later terse papers can still blur the worked route is the challenge-coverage family itself.
The maintained bundle already carries one exact piece-local spine from audience entry to one-piece answer, but those keys live across seven imported objects rather than one omnibus ledger.
\Cref{tab:workedexample-coverage-to-capsule-spine} writes that spine once for two contrastive cases: a surface-carried continuity half for the release-note audience, and a support-only recomputation hook for the auditor audience.
The point is not to create a new owner.
It is to keep visible that Synthesis~39--45 really travel through different maintained objects even when both routes start from one short emitted sentence.\cite{series:challengeroutes,series:challengestops,series:challengebranches,series:challengecoverage,series:challengesurfacehooks,series:challengeterminalwitnesses,series:challengecoveragewitnesspacks,series:challengecoveragewitnessslices,series:challengecoveragewitnesscores,series:challengeanswercapsules,series:workedexample}

\begin{table}[t]
\centering
\small
\begin{tabularx}{\linewidth}{@{}Y Y Y@{}}
\toprule
Worked case & Maintained piece-local spine & Why the split matters \\
\midrule
Release-note continuity piece & \nolinkurl{release_note_public_status_sentence} $\to$ \nolinkurl{public_status_sentence} $\to$ \nolinkurl{release_note_public_status_sentence_continuity_piece} $\to$ \nolinkurl{public_status_sentence_continuity_piece_coverage} $\to$ \nolinkurl{public_status_sentence_continuity_piece_surface} $\to$ \nolinkurl{public_status_sentence_continuity_piece_terminal} $\to$ \nolinkurl{public_status_sentence_continuity_piece_witness_core} $\to$ \nolinkurl{public_status_sentence_continuity_piece_witness_pack} $\to$ \nolinkurl{public_status_sentence_continuity_piece_coverage_release_note_slice} $\to$ \nolinkurl{public_status_sentence_continuity_piece_coverage_release_note_answer} & Surface-carried piece: the public wrapper sentence really does carry this half, so the spine must pass through exact surface realization before it reaches terminal ownership, proving factorization, and the final one-piece answer. \\
Auditor recomputation piece & \nolinkurl{auditor_compact_diff} $\to$ \nolinkurl{compact_diff_summary} $\to$ \nolinkurl{auditor_compact_diff_recompute_piece} $\to$ \nolinkurl{compact_numeric_diff_recompute_piece_coverage} $\to$ no surfaced hook $\to$ \nolinkurl{compact_numeric_diff_recompute_piece_terminal} $\to$ \nolinkurl{compact_numeric_diff_recompute_piece_witness_core} $\to$ \nolinkurl{compact_numeric_diff_recompute_piece_witness_pack} $\to$ \nolinkurl{compact_numeric_diff_recompute_piece_coverage_auditor_slice} $\to$ \nolinkurl{compact_numeric_diff_recompute_piece_coverage_auditor_answer} & Support-only piece: the compact diff sentence certifies the right to escalate, not the recomputation witness itself, so the spine intentionally skips surface realization while keeping exact terminal ownership and the audience-local proving delta visible. \\
\bottomrule
\end{tabularx}
\caption{A concrete worked coverage-to-capsule spine. The first row is the shortest surface-carried challenge route that still preserves exact realization, exact terminal ownership, shared-vs-audience-local proving burdens, and one-piece answer assembly. The second row is the parallel support-only route where no surfaced hook may be invented merely because later proving and answer objects still exist.}
\label{tab:workedexample-coverage-to-capsule-spine}
\end{table}

\paragraph{Tiny coverage-to-capsule drill.}
If a later terse paper asks whether the public wrapper still certifies continuity for a release-note reader, the honest worked path may pass through a surfaced hook and stop at the maintained one-piece continuity capsule.
If the paper instead asks whether the compact diff still justifies auditor recomputation, the honest route is support-only and must not invent a surfaced hook that the sentence never had.
That one contrast is enough to keep Synthesis~39--45 auditable in a brief paper without restating all seven companion notes inline.\cite{series:workedexample}


\begin{table}[t]
\centering
\small
\begin{tabularx}{\linewidth}{@{}Y Y Y@{}}
\toprule
Late-stage bucket & Canonical note range & First question to stop on there \\
\midrule
Release-local suffix & Synthesis~31--33 & Which maintained release-stage object owns the next successor claim? \\
Challenge routing / witnesses & Synthesis~36--44 & How does one contested semantic piece reach an approved terminal field? \\
Answer bundles / review handoffs & Synthesis~45--54 & What is the smallest sufficient maintained handoff for one audience/question pair? \\
Within-release source-only request layer & Synthesis~56--63 & What exactly was promised to readers, and what visible status/notice was emitted? \\
Response-packet reuse layer & Synthesis~64--67 & Which bounded handoff slice should travel, and how may it be reused? \\
Visible successor-refresh layer & Synthesis~68--70 & What exact reissue sentence is allowed once a reused packet survives the successor cut? \\
Successor-facing refresh-response closure & Synthesis~71--73 & Which successor-facing packet/menu cut closes the refreshed visible branch once one concrete follow-up is chosen? \\
\bottomrule
\end{tabularx}
\caption{Late-stage ownership in the worked bundle, grouped by task rather than restated as a long serial note list.}
\label{tab:workedexample-latestage-buckets}
\end{table}
The maintained cut now makes that claim concrete by storing a tiny \texttt{resolution\_contract} plus two typed \texttt{resolution\_examples} inside \texttt{example\_successor\_normal\_form.json}.
The contract fixes the step order \texttt{emit $\to$ clause $\to$ quote $\to$ cite $\to$ support $\to$ narrow} and requires every reusable sentence to admit at least one leftward trace to a narrower maintenance field. In parallel, the new \texttt{example\_release\_stage\_walkthrough.json} records the stage-by-stage release walk for the same successor pair so a reviewer can audit not only where a sentence stops, but also which stage in the larger maintenance spine owns the field family being compressed. The maintained cut now also ships \texttt{example\_successor\_challenge\_branches.json}, which turns the remaining implicit bridge into a concrete object: once one semantic piece of a route is actually challenged, the branch map states the exact escalation walk to the approved terminal field for that piece. It also ships \texttt{example\_successor\_challenge\_terminal\_witnesses.json}, which keeps exact piece-hook terminal ownership separate from both the sentence-family coverage certificate and the piece-to-segment surface-hook ledger. And the maintained cut now also ships \texttt{example\_successor\_challenge\_answer\_audit\_trails.json}, which starts one level later at the exact shipped answer carrier slot and then spells out the full leftward walk through the emitted answer card, whole-question answer sheet, one-piece capsules, piece-specific branches, and narrow terminal owners for that audience/question pair.
Each shipped example now also carries an explicit \texttt{terminal\_targets} list annotated with covered claim classes, a small \texttt{segments} list that binds auditable surface spans directly to those terminal targets, an ordered \texttt{surface\_assembly} list whose concatenation reconstructs the emitted sentence exactly from those segments plus literal glue text, and a tiny \texttt{comparison\_anchor} block that locks the sentence to this exact base$\to$successor pair plus the narrow fields and shipped object digests that make it true. Each anchor now also ships a compact \texttt{witness\_capsule} saying which already-declared provenance fields may travel beside unchanged wording, a smaller \texttt{witness\_floor} saying which minimum visible provenance fields should still appear whenever any such capsule is printed at all, an explicit \texttt{witness\_omittable} tail saying which optional visible-provenance fields may be dropped when a later note prints only that minimum bundle, and a tiny \texttt{inline\_provenance\_profiles} list naming the three allowed inline reuse modes as \texttt{silent}, \texttt{floor}, and \texttt{capsule}, a declared \texttt{default\_inline\_profile} fixing the archive-default reuse mode at \texttt{floor}, plus \texttt{recheck\_if}, \texttt{regenerate\_if}, and \texttt{refresh\_in\_place} lists separating object-cut drift that forces revalidation from pair/terminal drift that forces new wording before reuse and stating exactly which witness fields may be refreshed after successful recheck while the checked wording remains frozen.
So a reviewer no longer has to infer the stopping fields, the semantic split, the punctuation/order that turns the split spans back into the public sentence, \emph{whether} the same wording has been silently transplanted from some other successor comparison, \emph{whether} visible provenance is intentionally absent versus reduced to the minimum declared bundle, \emph{which} inline provenance mode the archive means a terse reuse site to inherit by default, \emph{which} tiny changes merely force recheck against a new object cut, \emph{which} witness fields may be refreshed in place after that recheck, \emph{which} visible provenance fields are mandatory versus omittable when space is tight, \emph{or} which changes retire that wording even if the template still reads plausibly.
For the outward-facing status sentence, the emitted clause resolves leftward through the clause pack, quote map, citation map, and support map, then stops at the continuity and closure tokens; the shipped segments split the surface text into the continuity half and the package-completeness half, the surface assembly records the semicolon and terminal period that join those halves back into the emitted sentence, the comparison anchor binds the compare/continuity/closure object digests so this sentence is only certified for the maintained worked-example object cut, the attached \texttt{witness\_capsule} says that only the release-pair ids, compare-report id, and digest witness may travel beside the unchanged wording as visible provenance, the attached \texttt{witness\_floor} says that whenever any such visible provenance is shown it should minimally name the base release, successor release, and compare-report cut, the attached \texttt{witness\_omittable} tail says that the digest witness may be dropped when a later note prints only that minimum inline bundle, and the attached \texttt{recheck\_if}/\texttt{regenerate\_if}/\texttt{refresh\_in\_place} lists now say that anchored-digest drift forces revalidation, that only \texttt{object\_digests} may be refreshed in place after successful recheck, and that a new release pair or changed continuity/closure terminal targets force reissue.
For the compact diff sentence, the trace is shorter: emitted sentence $\to$ clause pack $\to$ quote map $\to$ citation map $\to$ \texttt{example\_compare\_report.json:observed\_diffs}, with the continuity verdict naming which line items changed; the shipped segments then separate the two fallback-delta spans from the total-summary span, the surface assembly fixes the leading label plus semicolon separators so changed-family identity is bound only where the sentence actually asserts it, the comparison anchor prevents those numeric deltas from being reused as if they described a different successor pair or a different shipped compare-report cut, the attached \texttt{witness\_capsule} says that only the release-pair ids, compare-report id, and digest witness may travel beside the unchanged wording as visible provenance, the attached \texttt{witness\_floor} says that whenever any such visible provenance is shown it should minimally name the base release, successor release, and compare-report cut, the attached \texttt{witness\_omittable} tail says that the digest witness may be dropped when a later note prints only that minimum inline bundle, and the attached \texttt{recheck\_if}/\texttt{regenerate\_if}/\texttt{refresh\_in\_place} lists say that compare/continuity digest drift forces revalidation, that only \texttt{object\_digests} may be refreshed in place after successful recheck, and that pair drift or change to the numeric/changed-family terminal targets retires the sentence.
That turns the suffix into a non-decorative audit example rather than only a list of adjunct names, while keeping the stop-set, span-binding, exact sentence reconstruction, release-pair lock, object-cut witness, visible-provenance capsule, visible-provenance floor, explicit visible-provenance omission boundary, explicit recheck-vs-regenerate conditions, and the in-place refresh boundary in one stable wrapper instead of scattering them across prose.

\begin{table}[t]
\centering
\small
\begin{tabularx}{\linewidth}{@{}Y Y Y@{}}
\toprule
Emitted family & Terminal field(s) & Covered claim classes / why the trace stops there \\
\midrule
Public status sentence & \texttt{example\_successor\_continuity\_verdict.json:continuity\_decision}; \texttt{example\_publication\_closure\_verdict.json:closure\_status} & Covers \texttt{continuity\_decision} and \texttt{closure\_status}; these are the concrete decision tokens the wrapper sentence compresses. \\
Compact numeric diff & \texttt{example\_compare\_report.json:observed\_diffs}; changed-line-item slice in \texttt{example\_successor\_continuity\_verdict.json} & Covers \texttt{numeric\_delta} and \texttt{changed\_family\_identity}; numeric deltas and changed-family identity must resolve to narrower compare/continuity fields, not stop at reusable wording. \\
\bottomrule
\end{tabularx}
\caption{Terminal fields for the two typed sentence traces shipped in the worked normal-form object; these are bound machine-readably as \texttt{terminal\_targets}, while the emitted surface spans are bound as \texttt{segments}, reassembled by \texttt{surface\_assembly}, and locked to the maintained base$\to$successor pair plus shipped object digests by \texttt{comparison\_anchor} so a reviewer can see exactly which part of each sentence each endpoint discharges, that the declared spans really reconstruct the public sentence, which release comparison that sentence belongs to, which compare/verdict object cut it audits, and which witness fields may be refreshed in place after a successful recheck.}
\label{tab:workedexample-trace-endpoints}
\end{table}

In this example, \emph{all} control-plane drift that changes the effective reader set or request shaping (endpoint sets, AutoConf expansion rules, environment overrides, protocol filters, retrieval-mode switches) is treated as \textbf{state-surface drift} and therefore must bump \texttt{state\_decl\_id} even if the budget arithmetic is unchanged.\cite{series:deployedsurfaces,series:planstateequiv}

\paragraph{Three tiny examples.}
\begin{enumerate}[leftmargin=*]
\item \textbf{Window drift:} keep label \texttt{tw.reader\_24h.q500.v1} but change $Q$ (or the reset rule). The digest $TW$ changes and downstream $Q$-scale reasoning must be revisited.\cite{series:threatwindows}
\item \textbf{Surface drift:} keep label \texttt{enf.fallback.cct.v1} but add a new Tier~2 field. The digest $E_i$ changes; a user diff catches that the exposure surface widened.\cite{series:traceifaces}
\item \textbf{Evidence drift:} keep the profiling label but change the audit projection (bucket widths) or the confidence knob. The \texttt{evidence\_id} digest changes and auditors replay via the \texttt{replay\_hook}.\cite{series:receiptitems,series:auditorreplay}
\end{enumerate}

\paragraph{Claim-sameness card.}
A later terse paper should not have to infer whether two receipts are ``the same claim'' from prose.
The worked bundle therefore now makes the smallest practical sameness test explicit: if any row in \Cref{tab:workedexample-sameness-card} changes, the old claim must be treated as a different auditable object even when a headline budget number stays fixed.

\begin{table}[t]
\centering
\small
\begin{tabularx}{\linewidth}{@{}Y Y Y@{}}
\toprule
What changed? & Id that must change & Why the old claim is no longer the same object \\
\midrule
Reset rule, horizon, observer cadence, or public history scope & \texttt{tw\_id} & The repetition semantics changed, so the same local budget no longer means the same risk horizon. \\
Tier membership, transcript fields, bucket semantics, or stream interpretation & \texttt{exposure\_nf\_id} & The adversary now sees a different projection, so the old exposure surface is not the same claim. \\
Endpoint set, request-shaping defaults, downgrade policy, or fallback budget knobs & \texttt{state\_decl\_id} & The effective contacted-reader set or feasible transcript family changed, even if the arithmetic did not. \\
Audit menu, bucket map, confidence spend, or replay recipe & \texttt{evidence\_id} / \texttt{replay\_hook} & The evidence witness changed, so the old verification path no longer certifies the new statement. \\
\bottomrule
\end{tabularx}
\caption{Receipt-level claim sameness in the worked example. Stable labels are not enough; the digest-bound ids are the actual audit surface.}
\label{tab:workedexample-sameness-card}
\end{table}

\paragraph{Request-shaping drift (filters and budgets).}
If the effective delegated-routing request adds, removes, or changes request-shaping defaults or budgets (protocol filters, router allow/deny policy, path scope, routing/retrieval timeouts that govern downgrade and fallback), treat it as \textbf{state-surface drift}: it changes \emph{who is contacted} and what termination-time behavior is feasible, and therefore can change fallback rate and exposure.
Bump \texttt{state\_decl\_id} and treat it as recertification-triggering.\cite{series:deployedsurfaces,series:planstateequiv}

\paragraph{Anonymous-DHT public-owner cutover card.}
The sameness test above answers when the old receipt is no longer the same object.
Later terse anonymous-DHT papers still need the next routing answer: \emph{which narrow public owner should now speak first?}
\Cref{tab:workedexample-anondht-cutover-card} mirrors the current Boss Fight / evaluator / claim-surface companions in worked-example language, so a reviewer can translate one concrete receipt drift into one first public stop without rewalking the whole archive.\cite{series:bossfightA,series:bossfightB,series:bossfightC,series:evalnotary,series:abomoinl,series:releaseproperty}

\begin{table}[t]
\centering
\small
\begin{tabularx}{\linewidth}{@{}Y Y Y@{}}
\toprule
Worked drift pattern & First public owner & Why stop there first / widen only if \\
\midrule
Reset rule, horizon row, or public-history scope changes while the deployed projection is otherwise the same & Boss Fight~A & The repeated-use accountant and published horizon/cap changed. Widen to Boss Fight~B only if the dial-sheet contract changed too. \\
Same horizon and protected projection, but the declared observation proxy, dummy-cover width, retry width, timeout profile, or other knob vector changes & Boss Fight~B & The anonymous-DHT dial-sheet contract changed while the question is still \emph{which knob surface was declared}. Widen to Operational~A/B only if runtime conformance rather than declared knob choice is now live. \\
Same public cap and same dial-sheet family, same exported fallback/selection/summary scalars, but a checked \texttt{plan\_id}/\texttt{plan\_spec\_id} row in \texttt{example\_verifier\_report.json} or its support-bundle resolution in \texttt{example\_support\_bundle\_map.json} changes & Boss Fight~C & The same numbers no longer name the same replay contract. Widen to Eval~1 only if the live issue is attempted-evaluation notarization, public randomness, or the no-gap attempt prefix rather than the verifier object itself. \\
Same replayable verifier bundle, but the evaluator-attempt ledger, retry/randomness binding, or attempt-prefix discipline changes & Eval~1 & The live issue is notarized evaluation discipline rather than the declared bundle itself. \\
Same evaluator record, but the public anonymity-claim object or release-lineage receipt changes & ABOM/OINL or Release~A & The live issue is the signed claim surface or published release lineage rather than the accountant, dial sheet, or verifier contract. \\
\bottomrule
\end{tabularx}
\caption{Worked translation from receipt drift to the first narrow public owner that should speak next for an anonymous-DHT paper. This is a routing card, not a new theorem: it mirrors the standalone Boss Fight / evaluator / release notes so later terse papers can stop at the first already-owned public cut.}
\label{tab:workedexample-anondht-cutover-card}
\end{table}

\paragraph{Fallback-vector semantic guard.}
The row \texttt{fallback\_summary\_bits\_per\_lookup=0.056345507023} is retained because its three-term observation-attenuation arithmetic is reproducible.  It is not a deployment MC-EQ certificate: the bundle supplies no destination-indexed reference kernel, support-containment witness, positive reference floor, or joint-channel composition evidence.  Its maintained metadata therefore marks \texttt{semantic\_scope=illustrative\_tiered\_observation\_attenuation\_not\_mceq}, \texttt{composition\_evidence\_status=assumption\_only\_no\_joint\_channel\_witness}, and \texttt{publication\_eligible=false}.  Any route card below that previously used MC-EQ as the active theorem owner must now treat it as a hold-only branch until a real source-bound packet exists.

\paragraph{One exact verifier-side trilogy drill.}
Suppose the fallback line item keeps the same endpoint family and replay posture but raises the threat-window row from \texttt{tw.reader\_24h.q500.v1} to a stricter row with a new reset budget: the first honest public stop is Boss Fight~A, because the repeated-use horizon itself changed even before any dial-sheet or verifier question is live.
If the horizon row stays fixed but the declared observation proxy moves from $p_{\mathrm{obs}}=0.50$ to $0.60$ or the dummy cover shrinks from $w=1{,}034{,}140$ to $w=900{,}000$, the first honest public stop is Boss Fight~B, because the dial-sheet contract changed even if the same reader-private projection is still being protected.
If the horizon row and dial sheet are both held fixed, and \texttt{example\_verifier\_report.json} still exports the same four recomputed rows \texttt{fallback\_contact\_bits\_per\_lookup=0.04653341}, \texttt{fallback\_summary\_bits\_per\_lookup=0.056345507023}, \texttt{selection\_tax\_bits=0.263034405834}, and \texttt{total\_summary\_bits\_per\_lookup=0.265904562929}, but the deployment no longer checks the same five maintained plan/spec rows \texttt{eval-primary-v1}/\texttt{sha256:ea286a351bbfd0fd44f353b9b608f4bb3fbe6d535dae49c3a2621ec15e85f7ad}, \texttt{eval-fallback-cct-v1}/\texttt{sha256:e9168dbec2cabd9115a7f3ea9aebe0a75ff910d2295e450a340716cc7e541193}, \texttt{eval-fallback-vector-v1}/\texttt{sha256:38fd8a9252730b1eca7bb12cdd552a00cf50c98204615d197b1436fc061d72b5}, \texttt{eval-pathselect-v1}/\texttt{sha256:a4f3e91c6a43486caaa6d31b3330927284369aaeed4e0e6cf6683272ac59b600}, and \texttt{eval-profiling-eq-v1}/\texttt{sha256:db245ce4f3c200322eea17a42a2c64cf9dc28bdeb18fa24e5d3643b51eb1e6f6}, or no longer resolves the same five maintained support bundles \texttt{bundle://worked-example/primary}, \texttt{bundle://worked-example/fallback-cct}, \texttt{bundle://worked-example/fallback-vector}, \texttt{bundle://worked-example/pathselect}, and \texttt{bundle://worked-example/profiling-eq} in \texttt{example\_support\_bundle\_map.json}, then the first honest public stop is Boss Fight~C: the old public card is stale even though the published numbers themselves have not moved.
That gives later terse papers one concrete worked translation from generic receipt drift to the first public Anonymous-DHT owner that already carries the right fail-closed cutover, while leaving evaluator-attempt notarization and release closure with their narrower downstream owners.\cite{series:bossfightA,series:bossfightB,series:bossfightC,series:evalnotary,series:releaseproperty}



\paragraph{Route-shape selection card.}
The first-owner cutover card says \emph{who} should speak next.
A later terse paper still needs one more answer before it cites anything: \emph{what public route shape is now live?}
The worked example now keeps that second choice explicit as well, so a familiar receipt drift is not silently flattened into the wrong kind of spine.\cite{series:bossfightA,series:bossfightB,series:bossfightC,series:mceq,series:abomoinl,series:releaseproperty,series:state1,series:state1A,series:mucc,series:state3,series:state3A,series:eval1A,series:eval2,series:eval3}

\begin{table}[t]
\centering
\small
\begin{tabularx}{\linewidth}{@{}Y Y Y@{}}
\toprule
Worked trigger after the first-owner cut & Route shape to cite next & Stop here first / widen only if \\
\midrule
Horizon, dial, or checker drift inside the same protected lookup family & Linear accountant $\to$ dial sheet $\to$ replay route (Boss Fight~A $\to$ Boss Fight~B $\to$ Boss Fight~C) & Stop at the first changed owner among the three. Widen only if the live question truly moves from repeated-use accountant, to declared knob surface, to replay contract. \\
Same theorem packet, but the digest-bound public claim object or the signed lineage receipt changes & Hold-only theorem $\to$ claim $\to$ release route (MC-EQ [hold] $\to$ ABOM/OINL $\to$ Release~A) & No current worked row may enter this route as MC-EQ evidence. Reopen MC-EQ only for a source-bound TV or support-disciplined card; then stop at ABOM/OINL for the exact claim object and at Release~A only for the signed publication / lineage verdict. \\
State-surface, committee-contact floor, or public monitoring/alarm drift & Branched state route (State~1 (+State~1A) / State~2 / State~3 (+State~3A)) & Stop at State~1 for the generic state accountant, at State~2 when the committee-contact specialization is live, and at State~3 when the public monitoring contract is live; keep State~1A and State~3A as optional support branches rather than compulsory middle stops. \\
Same public claim family, but evaluator-attempt discipline, audit-transfer contract, or cadence-facing planning sentence changes & Evaluation route (Eval~1 (+Eval~1A) $\to$ Eval~2 $\to$ Eval~3) & Stop at Eval~1 for the retry-safe notarized evaluation object, widen to Eval~2 only when the live issue is audit/live transfer, and widen to Eval~3 only for cadence-facing conversion or planning sentences. \\
\bottomrule
\end{tabularx}
\caption{After identifying the first public owner, the worked example still classifies the route shape that should carry the next citation. This keeps later terse papers from flattening claim objects, monitoring branches, or evaluator planning packets into the wrong spine.}
\label{tab:workedexample-route-shape-card}
\end{table}

\paragraph{One tiny route-shape drill.}
If the receipt drift is only a new checker id under the same Boss Fight inequality, the honest next citation is still on the linear accountant $\to$ dial $\to$ replay route, and the first stop is Boss Fight~C rather than ABOM/OINL or Release~A.\cite{series:bossfightA,series:bossfightB,series:bossfightC}
The MC-EQ branch is currently hold-only. A future source-bound MC-EQ theorem packet may re-enter the theorem $\to$ claim $\to$ release route, but the illustrative fallback-vector row cannot trigger that route. Once a valid theorem packet exists, a changed digest-bound claim object or signed lineage receipt belongs to ABOM/OINL or Release~A rather than to the worked receipt.\cite{series:mceq,series:abomoinl,series:releaseproperty}
If the same receipt now raises a committee-contact floor question together with a public alarm question, the honest route shape is branched: State~1 at the root, State~2 for the anonymous-DHT specialization, and State~3 for the monitoring contract, with State~1A / State~3A staying optional support packets.\cite{series:state1,series:state1A,series:mucc,series:state3,series:state3A}
That extra route-shape check keeps the public-owner cutover card from becoming another flattening shortcut.


\paragraph{One tiny notary packet drill.}
The maintained worked bundle now carries the retry-safe evaluation card directly as structured verifier data rather than prose-only memory. In \nolinkurl{example_verifier_report.json}, \nolinkurl{notary_replay_packet} fixes receipt digest \nolinkurl{095e828fd09ef186757f5dd49e84769603a2b59ea77c4a8bfb0edb720bc64cc1}, release id \nolinkurl{worked-example-draft-199}, artifact note version \nolinkurl{1.99}, the fallback-contact line item, replay hook \nolinkurl{eval-fallback-cct-v1}/\nolinkurl{sha256:e9168dbec2cabd9115a7f3ea9aebe0a75ff910d2295e450a340716cc7e541193}, logged prefix \nolinkurl{attempts_present=[1,2]} with \nolinkurl{first_passing_attempt=2}, the disclosed retry budget \nolinkurl{1e-6 total across sched.delta.4step.v1}, the five-metric Bonferroni rule with threshold \nolinkurl{6e-8}, the verifier line \nolinkurl{reported_p_value=4e-8}, the lower replay support cut \nolinkurl{bundle://worked-example/fallback-cct}, the optional paired support packet \nolinkurl{example_uvi.json}, the release-side handoff role \nolinkurl{fresh_replay_verdict}, and the quotable output \nolinkurl{fallback_contact_bits_per_lookup=0.04653341}. So later terse papers may now stop at one exact worked packet when the live issue is the retry-safe notary object itself, widen to Eval~1A only when checkpoint / witness-policy hardening is itself live, widen to Eval~2 only when the live issue becomes audit/live transfer, widen to Eval~3 only for cadence-facing conversion or planning, and widen to Release~A only once the live issue becomes closure packaging above that same replay verdict.\cite{series:evalnotary,series:eval1A,series:advantagecontracts,series:eval3,series:releaseproperty,series:workedexample}

\paragraph{One calibration owner-map boundary drill.}
The current worked-example cut exposes the Eval~3 calibration route through \nolinkurl{example\_line\_item\_owner\_map.json} and the generated series spine rather than through a new verifier/support packet. The owner map includes \nolinkurl{calibration\_recipe\_packet} and keeps State~1 / State~2 / State~3 / MC-EQ as imported theorem/accountant and destination-equalization neighbors while Eval~3 owns the local numbers-first conversion stop. However, \nolinkurl{example\_verifier\_report.json} does not contain \nolinkurl{calibration\_recipe\_packet}; \nolinkurl{example\_replay\_plans.json} does not contain \nolinkurl{eval-calibration-recipe-v1}; and \nolinkurl{example\_support\_bundle\_map.json} does not contain \nolinkurl{bundle://worked-example/calibration-recipe}. The honest drill is therefore a boundary drill, not a packet drill: stop at Eval~3 for the local horizon/cadence/alert-tax conversion object, reopen State~1 / State~2 / State~3 or MC-EQ when the imported theorem/accountant or destination row becomes live, and widen to Eval~1 / Eval~2 / Release~A only after the live issue has left the conversion layer.\cite{state1,state2,state3,series:mceq,series:eval3,series:evalnotary,series:advantagecontracts,series:releaseproperty,series:workedexample,series:seriesspines}

\paragraph{One monitoring-contract absence boundary drill.}
The current worked-example cut does not carry the public monitoring card as a separate packet. \nolinkurl{example_state_decl.json} does not contain \nolinkurl{monitoring_contract_packet}; \nolinkurl{example_replay_plans.json} does not contain \nolinkurl{eval-cppc-v1}; and \nolinkurl{example_support_bundle_map.json} does not resolve \nolinkurl{bundle://worked-example/cppc}. So later terse papers should stop at State~3 for the local CPPC card when the live issue is the public monitoring contract itself, widen to State~3A only for the optional schema/linter/many-detector support packet, and widen farther to Eval~3 / Eval~1 / Operational~B / Release~A only once the live issue has already left the monitoring contract.\cite{series:state3,series:state3A,series:eval3,series:evalnotary,series:opB,series:release}

\section{Release evolution as a stable spine}
The worked bundle now treats successor maintenance as a nine-stage chain:
\emph{diff $\to$ classify $\to$ obligate $\to$ close $\to$ summarize $\to$ support $\to$ cite $\to$ quote $\to$ clause}.
This is not a new receipt primitive.
It is a compact answer to a recurring publication problem: later notes often need to say \emph{how} a released claim evolved without repeating the compare/replay/closure story in prose.
The concrete dependency DAG is still carried by \texttt{example\_successor\_derivation\_graph.json}; the maintained \texttt{example\_release\_spine.json} now freezes the stable stage vocabulary, the tiny per-stage contracts, and an explicit release-local stop-here / widen-only-if rule for every stage-order row, while the companion \texttt{example\_release\_stage\_walkthrough.json} shows the same successor comparison stage-by-stage with concrete artifacts, owned field families, resolved inputs, and citation boundaries. The worked bundle now also makes the obligate$\to$close bridge explicit with \texttt{example\_release\_closure\_ledger.json}, so a reviewer can inspect role-by-role package satisfaction before reading the compressed closure token. That turns the spine from a decorative pipeline diagram into a checkable release-evolution interface. Synthesis~32 owns that stage interface and Synthesis~33 owns the closure bridge; this worked note just ships the concrete instantiation.\cite{series:releasespine,series:releaseclosure} The nine-stage release spine plugs into the six-stop stable ladder at the release anchor: downstream notes can cite one maintained release-local stage cut and then widen to the disclosure-packet and review owners only if the question actually needs those later stops.

\paragraph{Earliest sufficient stage cut (worked rule).}
If the question is only ``what changed?'' the worked path stops at \texttt{example\_compare\_report.json}. If it is ``what had to be replayed or re-certified?'' it stops at the continuity verdict. If it is ``what must be republished?'' it stops at the obligation profile. If it is ``was the successor package actually complete?'' it must reach the closure ledger and closure verdict. Only when a later note needs emitted public wording should it continue to the summarize/support/cite/quote/clause suffix. Machine-readably, those release-local cuts now sit in the release spine itself as stage-order rows with explicit \texttt{stop\_when} and \texttt{widen\_only\_if} clauses, so the stable series spine in Synthesis~55 can attach one release-local cut to the theorem root on the left and the referee handoff on the right without letting this worked note absorb those neighboring owners.\cite{series:releasespine,series:seriesspines}

\paragraph{Stage-to-stop interlock.}
The nine release stages stay honest only because they collapse back onto the six first-sufficient audit stops rather than pretending to be a second independent ladder.
The maintained bundle therefore exposes one explicit interlock: diff/classify land on receipt/replay owners; obligate/close land on the release-local package-completeness cut; summarize lands on the narrow public-status wrapper; support/cite/quote/clause land on the answer-side archive cut and review packaging; and only a genuine crossing of that archive cut widens into the adjacent source-only request ladder.
That is the concrete worked interface Synthesis~55 reuses when it treats release evolution as a stable middle bridge instead of a new public story.\cite{series:releasespine,series:releaseclosure,series:statusenvelopes,series:seriesspines,series:challengeanswerdisclosurepackets,series:challengeanswerreviewmenus,series:pairedterminalcutoverrules}

\begin{table}[t]
\centering
\small
\begin{tabularx}{\linewidth}{@{}Y Y Y Y@{}}
\toprule
Stage & Canonical artifact & Owned output family & Question answered \\
\midrule
Diff & \texttt{example\_compare\_report.json} & \texttt{observed\_diffs}, \texttt{classification} & What changed between base and successor? \\
Classify & \texttt{example\_successor\_continuity\_verdict.json} & \texttt{continuity\_decision}, \texttt{changed\_line\_items}, \texttt{replay\_plan\_ids} & Which replay scope and continuity class follow from that diff? \\
Obligate & \texttt{example\_release\_obligation\_profile.json} (+ action matrix) & \texttt{classes}, \texttt{required\_outputs} & Which public outputs must now be refreshed? \\
Close & \texttt{example\_publication\_closure\_verdict.json} & \texttt{closure\_status}, \texttt{provided\_publication\_outputs} & Were the required outputs actually shipped? \\
Summarize & \texttt{example\_successor\_lineage\_notice.json} & \texttt{status\_summary} & What is the outward-facing successor status sentence? \\
Support & \texttt{example\_successor\_claim\_support\_map.json} & \texttt{claims} & Which narrower object supports each public clause? \\
Cite & \texttt{example\_successor\_citation\_map.json} & \texttt{queries} & What is the smallest sufficient citation target downstream? \\
Quote & \texttt{example\_successor\_quote\_map.json} & \texttt{quote\_targets} & Which exact field should later notes quote? \\
Clause & \texttt{example\_successor\_clause\_pack.json} & \texttt{clauses} & Which short reusable clause is already checked? \\
\bottomrule
\end{tabularx}
\caption{The worked example's release-evolution spine. The derivation graph owns the concrete dependency DAG; the spine now owns both the stable stage vocabulary and the tiny per-stage handoff contract saying which field family each stage may export downstream.}
\label{tab:release-spine}
\end{table}

\begin{table}[t]
\centering
\small
\begin{tabularx}{\linewidth}{@{}Y Y Y@{}}
\toprule
Release-stage slice & First sufficient audit stop & Widen only if \\
\midrule
Diff $\to$ Classify & Receipt / replay stop (compare report plus continuity verdict) & The live issue becomes required public outputs or package completeness rather than replay scope. \\
Obligate $\to$ Close & Release-local package-completeness cut (obligation profile plus closure ledger / verdict) & The live issue becomes emitted public wording, downstream support, or checked-answer handoff. \\
Summarize & Narrow public-status wrapper (lineage notice and status token) & The live issue becomes the exact downstream support/citation/quote/clause owner rather than the wrapper sentence. \\
Support $\to$ Cite $\to$ Quote $\to$ Clause & Answer-side archive cut and review packaging & The live issue has genuinely crossed the answer-side archive cut into omitted-support routing. \\
After the archive cut is crossed & Request-side terminal form for the adjacent source-only branch & Any named reopen trigger fires, so terminal shortcuts are stale and the paper must fail closed to the reopening owner. \\
\bottomrule
\end{tabularx}
\caption{Worked stage-to-stop interlock. The release spine is auditable precisely because each stage slice collapses onto one first-sufficient stop and names the next honest widening boundary, rather than becoming a decorative second ladder.}
\label{tab:workedexample-stage-stop-interlock}
\end{table}

\paragraph{Tiny release-local clause drill.}
Suppose a later terse paper needs only the sentence 
\emph{``the successor diff stayed within the same checked public-status wrapper and the refreshed package already closed.''}
The honest worked route is:
\nolinkurl{example_compare_report.json}
$\to$
\texttt{example\_successor\_continuity\_verdict.json}
$\to$
\texttt{example\_release\_obligation\_profile.json}
$\to$
\texttt{example\_publication\_closure\_verdict.json}
$\to$
\texttt{example\_successor\_lineage\_notice.json}.
That route is enough because the live issue is still release-local: what changed, what had to be refreshed, whether the package closed, and which narrow wrapper sentence is now honest. The paper should widen to the answer-side terminal form only when the live question becomes the exact checked public answer behind that wrapper, and should cross the archive cut to the bounded request-side packet only when omitted support behind that checked answer is actually being requested. This is the concrete worked counterpart of the abstract stable-spine release-slot discipline in Synthesis~55.\cite{series:releasespine,series:releaseclosure,series:statusenvelopes,series:seriesspines,series:pairedterminalcutoverrules}

\begin{table}[t]
\centering
\small
\begin{tabularx}{\linewidth}{@{}Y Y Y@{}}
\toprule
Question family & First sufficient owner & Why stop there \\
\midrule
Theorem or lower-bound justification & Mathematical root via \texttt{example\_line\_item\_owner\_map.json} / Synthesis~24 & No receipt, replay, or review wrapper is needed if the task is only mechanism justification. \\
Public claim identification & \texttt{example\_receipt.json} & The receipt owns the released claim surface and should answer identification questions without widening into replay or maintenance prose. \\
Replay / recomputation & \texttt{example\_replay\_plans.json} & Replay questions stop at the plan catalog and evidence hooks rather than drifting into successor wording. \\
Release evolution / package completeness & \texttt{example\_release\_spine.json} plus continuity / closure verdicts & This is the first owner that says what moved, what had to be refreshed, and whether the successor package closed. \\
Answer-side archive cut for omitted support & \texttt{example\_successor\_challenge\_answer\_disclosure\_packets.json} & The disclosure packet is the first sufficient answer-side owner for the exact requestable archive slice behind one shipped public answer. \\
Referee follow-up packaging & \texttt{example\_successor\_challenge\_answer\_review\_menus.json} & Once the task widens past that archive cut into one concrete follow-up class, the review menu is the smallest sufficient owner. \\
Paper-level source-only promise / visible notice & Source-only owner bucket Synthesis~56--63 & Contract, completion/status/carryforward, and notice-family choice are adjacent source-only owners, not extra spine segments. \\
Bounded response packet / reuse rules & Source-only owner bucket Synthesis~64--67 & This is the first bucket that owns exact packet cuts and their carryforward/delta rules. \\
Visible refresh notice / family choice & Source-only owner bucket Synthesis~68--70 & Visible successor-refresh wording has its own owner bucket before any successor-facing follow-up choice is made. \\
Successor-facing refresh-response packet / closure & Source-only owner bucket Synthesis~71--73 & This final bucket owns the successor-facing menu choice, exact packet cut, and capstone closure verdict. \\
\bottomrule
\end{tabularx}
\caption{Human summary of the worked example's machine-readable first-sufficient-owner routing table; the maintained \texttt{example\_question\_routes.json} adjunct binds these stops to concrete question stems and widening transitions.}
\label{tab:workedexample-routingcuts}
\end{table}

The maintained \texttt{example\_question\_routes.json} adjunct is deliberately tiny. It does not create a new owner or a new review layer. It turns the six-step stable ladder into a machine-readable routing catalog and mirrors the adjacent eighteen-step source-only request ladder as well: each generic task gets one first sufficient stop, each source-only request subtask gets one explicit rung with its own stop-here / widen-only-if rule, and each shipped audience/question answer binds one public anchor, one emitted answer / sentence-lock pair, one inherited brief/expanded/exact/full-audit publication-budget map, one mirrored answer-review request-class / answer-escalation-step / release-stage vocabulary posture bundle, one disclosure packet with its support-surface / manifest-cut / inventory-cut basis, one notice-keyed omitted-support packet cut with its followup-vocabulary / scope / family-basis / support-surface / manifest-cut / inventory-cut mirrors, one visible-refresh successor packet cut with its family-basis / support-surface / manifest-cut / inventory-cut / validation / closure mirrors, one review menu, and only the adjacent source-only owner buckets that matter if the question widens beyond the answer-side archive cut. That keeps the worked example auditable without forcing later terse papers to restate the whole source-only maintenance tail or to guess which source-only owner is first sufficient once the question has left the answer-side archive cut.

\section{One concrete audience-question drill}
\label{sec:drill}
To keep the worked example non-decorative, this note fixes one complete walk for the shipped auditor question \emph{``what changed in the compact diff?''}
The walk is not a new owner.
It is merely the smallest public-to-review extraction already implied by \texttt{example\_question\_routes.json}, \texttt{example\_successor\_challenge\_answer\_audit\_trails.json}, \texttt{example\_successor\_challenge\_answer\_disclosure\_packets.json}, and \texttt{example\_successor\_challenge\_answer\_review\_menus.json}.
That extraction is the concrete shape later terse papers should cite instead of paraphrasing a large tail.

\begin{table}[t]
\centering
\small
\begin{tabularx}{\linewidth}{@{}Y Y Y Y@{}}
\toprule
Step & Artifact / key & Settles & Widen only if \\
\midrule
Public anchor & \texttt{example\_compare\_report.json} & The emitted public compact-diff answer really belongs to the maintained base$\to$successor comparison. & Widen only if the reader needs checked wording rather than mere claim identification. \\
Checked wording & \texttt{auditor\_compact\_diff\_summary\_answer\_card} with sentence lock \texttt{compact\_numeric\_diff} & The exact public sentence and its reuse posture are fixed. & Widen only if the reader needs the exact shipped clause location or a leftward audit walk. \\
Exact carrier & Citation docket $\to$ carrier slot \texttt{auditor\_compact\_diff\_summary\_answer\_carrier\_slot} & The exact shipped clause carrying the sentence is fixed without reopening the whole answer family. & Widen only if the reader challenges one semantic piece of the sentence. \\
Leftward audit & Audit trail \texttt{auditor\_compact\_diff\_summary\_answer\_audit\_trail} & The sentence decomposes into named semantic pieces with approved terminal owners. & Widen only if the reader needs the requestable archive slice or a referee-selected follow-up class. \\
Answer-side archive cut & Disclosure packet \texttt{auditor\_compact\_diff\_summary\_answer\_disclosure\_packet} & The exact request packet, evidence classes, validation companions, and answer-side omitted-support cut are fixed. & Widen only if one concrete referee handoff class must now be chosen. \\
Referee handoff & Review menu \texttt{auditor\_compact\_diff\_summary\_answer\_review\_menu} & The six smallest-sufficient follow-up classes are fixed: brief, expanded, exact-location, full-audit, request-packet, or reveal-order. & Stop here unless the chosen class itself reopens a later source-only owner bucket. \\
\bottomrule
\end{tabularx}
\caption{A compact audited walk for one shipped audience/question answer. This is the worked example in its most reusable form: one public anchor, one checked sentence, one exact carrier, one leftward audit trail, one answer-side archive cut, and one review menu.}
\label{tab:workedexample-concrete-drill}
\end{table}

The maintained route object also fixes the exact reusable tuple, not just the six prose headings:
\texttt{example\_compare\_report.json} $\to$ answer card \texttt{auditor\_compact\_diff\_summary\_answer\_card} $\to$ carrier slot \texttt{auditor\_compact\_diff\_summary\_answer\_carrier\_slot} $\to$ audit trail \texttt{auditor\_compact\_diff\_summary\_answer\_audit\_trail} $\to$ disclosure packet \texttt{auditor\_compact\_diff\_summary\_answer\_disclosure\_packet} $\to$ review menu \texttt{auditor\_compact\_diff\_summary\_answer\_review\_menu}.
Its default inline publication profile is \texttt{floor}, and the shipped visible provenance floor is only \texttt{base\_release\_id}, \texttt{successor\_release\_id}, and \texttt{compare\_report\_id}. For later terse papers, those same answer-side owners now also compress to one reusable four-anchor answer-side terminal citation normal form: \nolinkurl{example_compare_report.json} $\to$ \texttt{auditor\_compact\_diff\_summary\_answer\_card} $\to$ \texttt{auditor\_compact\_diff\_summary\_answer\_disclosure\_packet} $\to$ \texttt{auditor\_compact\_diff\_summary\_answer\_review\_menu}. The imported publication profile, citation docket, carrier slot, audit trail, and escalation ladder remain auditable support for that shorter normal form and reopen only when exact-location, full-audit, or reveal-order is itself the live question. The worked review-menu ledger now ships that shorter stop explicitly as \nolinkurl{answer_side_terminal_citation_normal_form} together with its one-clause \nolinkurl{answer_side_terminal_sentence_normal_form}, and \nolinkurl{example_question_routes.json} plus \nolinkurl{example_series_spine.json} mirror both fields verbatim so later terse papers can reuse the maintained stop without reassembling it from the longer answer tuple. Those same mirrors now also carry \nolinkurl{answer_side_terminal_stop_when}, \nolinkurl{answer_side_terminal_widen_only_if}, \nolinkurl{answer_side_terminal_reopen_artifact}, and \nolinkurl{answer_side_terminal_first_reopen_owner_map}, so the answer-side shortcut itself is now a self-describing worked cut rather than a prose-only prelude to the paired-terminal papers.
The first adjacent source-only widening step beneath that exact answer-side tuple is also now explicit rather than merely nearby: \nolinkurl{example_public_request_notice.json} $\to$ \texttt{worked\_example\_receipt\_interlock\_successor\_public\_request\_response\_menu} $\to$ \texttt{worked\_example\_receipt\_interlock\_successor\_public\_request\_response\_menu\_\_request\_scope\_check}, and that bounded packet carries only \nolinkurl{../tools/validate_example.py}, \nolinkurl{example_artifact_inventory.json}, \nolinkurl{example_public_request_contracts.json}, \nolinkurl{example_validation_report.json}, and \nolinkurl{support_manifest.json}. The next visible-refresh continuation is now exact as well rather than generic: \texttt{worked\_example\_receipt\_interlock\_successor\_public\_request\_response\_menu\_\_request\_scope\_check\_\_delta\_\_refresh\_notice\_\_selection\_\_response\_menu} $\to$ \texttt{worked\_example\_receipt\_interlock\_successor\_public\_request\_response\_menu\_\_request\_scope\_check\_\_delta\_\_refresh\_notice\_\_selection\_\_response\_menu\_\_packet\_membership\_check} $\to$ \texttt{worked\_example\_receipt\_interlock\_successor\_public\_request\_response\_menu\_\_request\_scope\_check\_\_delta\_\_refresh\_notice\_\_selection\_\_response\_menu\_\_packet\_membership\_check\_\_closure}; that successor-facing packet carries only \nolinkurl{example_public_request_response_packets.json} and \nolinkurl{example_public_request_response_packet_refresh_notices.json}, and the closure verdict is \texttt{closed\_finite\_basis} over \texttt{public\_paths\_only} with reopen triggers \texttt{followup\_taxonomy\_expands}, \texttt{packet\_reuse\_reopens}, and \texttt{visible\_refresh\_sentence\_rebinds\_to\_new\_delta\_family}. For later terse papers, those four maintained anchors now serve as one reusable request-side terminal citation normal form: opening notice, first bounded omitted-support packet, exact successor-facing visible-refresh packet, and exact closure verdict. Taken together, the answer-side stop and the adjacent request-side capstone instantiate the paired terminal citation discipline of Synthesis~74, while Synthesis~75 makes the cutover rule explicit: stop at the answer-side terminal normal form unless the live question has become exact-location, full-audit, or reveal-order, or has genuinely crossed the answer-side archive cut into the adjacent source-only branch; if it has crossed, widen only while no named request-side reopen trigger has fired, and otherwise fail closed to the named reopening owner.\cite{series:pairedterminalcitationdiscipline,series:pairedterminalcutoverrules}
That is the practical normal form later terse papers should cite: the compared release pair, the maintained diff owner, the exact checked sentence, then one exact source-only contract packet if a referee widens past the answer-side review menu, and then one exact visible-refresh packet/closure pair if that widening continues onto the successor-facing branch.

A practical cutover card now falls straight out of that maintained route.
\begin{enumerate}[leftmargin=1.5em]
\item If the live question is only \emph{``is this the right shipped compact-diff sentence?''}, stop at the answer-side terminal form. It already fixes the compared public anchor, the checked sentence, the answer-side archive cut, and the referee handoff.
\item If the live question becomes \emph{exact-location}, \emph{full-audit}, or \emph{reveal-order}, reopen the imported answer-side middle owner instead. That is exactly when the citation docket, carrier slot, audit trail, or escalation ladder becomes live again.
\item If the live question has crossed the answer-side archive cut but is still only \emph{``what is the first source-only packet just beyond that cut?''}, stop at \nolinkurl{example_public_request_notice.json} $\to$ \texttt{worked\_example\_receipt\_interlock\_successor\_public\_request\_response\_menu\_\_request\_scope\_check}. That bounded omitted-support packet is the first honest request-side stop before any visible-refresh question is live.
\item If the live question has widened further to \emph{``where does the current visible-refresh branch close?''}, the request-side terminal form is the honest capstone. It preserves the opening notice, first bounded packet, exact successor-facing packet, and exact closure verdict for the same branch.
\item If \texttt{followup\_taxonomy\_expands}, \texttt{packet\_reuse\_reopens}, or \texttt{visible\_refresh\_sentence\_rebinds\_to\_new\_delta\_family} fires, fail closed and reopen the named request-side owner. A terse paper should not cite any terminal shortcut once that branch has gone stale.
\end{enumerate}
\noindent This card makes the worked example usable under real space pressure: later terse papers can see, question by question, when the answer-side terminal form is enough, when the first bounded source-only packet is the honest next stop, when the request-side capstone is warranted, and when they must fail closed to a reopened owner rather than citing any shortcut.


\paragraph{Three-check referee recipe.}
A referee who wants to audit the worked route without reopening the whole tail can do it in three machine-checkable yes/no checks.
The point is not to create a new owner.
It is to expose the smallest auditable check sequence already implicit in the shipped compare walk, route catalog, and closure verdict ledger.
\begin{enumerate}[leftmargin=1.5em]
\item \textbf{Compare-layer identity.} In \nolinkurl{example\_compare\_walkthrough.json}, the worked comparison should still land in \nolinkurl{classification\_result=replay\_changed\_surfaces} with \nolinkurl{continuity\_decision=same\_certified\_interface\_replayed\_surfaces}, \nolinkurl{changed\_line\_items=[fallback\_contact\_surface,fallback\_tiered\_vector]}, and \nolinkurl{replay\_plan\_ids=[eval\_fallback\_cct\_v1,eval\_fallback\_vector\_v1]}. If not, reopen the compare report, continuity verdict, and release-obligation profile instead of treating the compact-diff sentence as self-authenticating.
\item \textbf{Answer-side terminal honesty.} In \nolinkurl{example\_question\_routes.json}, the maintained \nolinkurl{answer\_side\_terminal\_citation\_normal\_form} should still have the four anchors \nolinkurl{public\_anchor $\to$ checked\_answer\_card $\to$ answer\_disclosure\_packet $\to$ answer\_review\_menu}, with \nolinkurl{reopen\_if\_request\_classes=[exact\_location\_check,full\_audit\_check,reveal\_order\_check]}. If not, reopen the citation-docket / carrier-slot / audit-trail / escalation-ladder owners, or the review-menu owner, as soon as one of those request classes becomes live.
\item \textbf{Request-side closure honesty.} In \nolinkurl{example\_public\_request\_response\_packet\_refresh\_response\_packet\_closure\_verdicts.json}, the closure key \nolinkurl{worked\_example\_receipt\_interlock\_successor\_public\_request\_response\_menu\_\_request\_scope\_check\_\_delta\_\_refresh\_notice\_\_selection\_\_response\_menu\_\_packet\_membership\_check\_\_closure} should still certify \nolinkurl{closure\_status=closed\_finite\_basis} over \nolinkurl{closure\_basis\_scope=public\_paths\_only} with \nolinkurl{reopen\_trigger\_tokens=[followup\_taxonomy\_expands,packet\_reuse\_reopens,visible\_refresh\_sentence\_rebinds\_to\_new\_delta\_family]}. If not, reopen exactly the owner named in \nolinkurl{reopen\_owner\_map}: the successor refresh-response menu for taxonomy growth, the packet-carryforward profile for reuse drift, or the refresh-notice selection for delta-family rebinding.
\end{enumerate}
\noindent These three checks are intentionally asymmetric.
The first guards sameness of the compared release pair, the second guards the paper-facing checked-answer shortcut, and the third guards the adjacent source-only visible-refresh capstone.
Together they give later terse papers one auditable recipe for deciding whether they may still stop at the paired terminal forms or must reopen a narrower owner.

\paragraph{Dual terminal normal-form card.}
Later terse papers often need the paired cutover in one glance rather than as another prose paragraph.
\Cref{tab:workedexample-dual-terminal-card} therefore places the two exact terminal forms and their paired shortcut side by side.
The answer-side form certifies the shipped checked sentence plus the answer-side archive cut for the compared release pair.
The request-side form certifies the first bounded omitted-support packet, the exact successor-facing visible-refresh packet, and the current closure verdict for that same branch.
The paired shortcut is honest only when the answer-first guard still holds, the question has genuinely crossed the answer-side archive cut, and no named request-side reopen trigger has fired.

\begin{table}[t]
\centering
\small
\begin{tabularx}{\linewidth}{@{}Y Y Y@{}}
\toprule
Form & Certifies exactly & Reopen / stale condition \\
\midrule
Answer-side terminal form & The checked compact-diff sentence for this compared release pair, the exact answer-side omitted-support cut behind that sentence, and the first referee handoff that may widen it. & Reopen when exact-location, full-audit, or reveal-order becomes the live question, or when the question has genuinely crossed the answer-side archive cut. \\
Request-side terminal form & The first bounded source-only packet beyond that archive cut, the exact successor-facing visible-refresh packet for the same branch, and the current closure verdict and closure scope for that branch. & Reopen when any named request-side reopen trigger fires. \\
Paired terminal shortcut & One compact paired reference saying both ``this is the checked shipped sentence'' and ``this is the adjacent source-only branch capstone'' without reopening the whole review spine. & Honest only while the answer-first guard and adjacency guard both hold; stale as soon as either terminal form is stale. \\
\bottomrule
\end{tabularx}
\caption{Two exact terminal forms plus the paired shortcut they support. This is the smallest referee-facing card for deciding whether a later terse paper should stop, widen once, or fail closed.}
\label{tab:workedexample-dual-terminal-card}
\end{table}

\begin{quote}\small
\textbf{Ordered anchors.}
The answer-side terminal form is
\path{example_compare_report.json}\\
$\to$ \texttt{auditor\_compact\_diff\_summary\_answer\_card}\\
$\to$ \texttt{auditor\_compact\_diff\_summary\_answer\_disclosure\_packet}\\
$\to$ \texttt{auditor\_compact\_diff\_summary\_answer\_review\_menu}.

\smallskip
\noindent The request-side terminal form is
\path{example_public_request_notice.json}\\
$\to$ \texttt{worked\_example\_receipt\_interlock\_successor\_public\_request\_response\_menu\_\_request\_scope\_check}\\
$\to$ \texttt{worked\_example\_receipt\_interlock\_successor\_public\_request\_response\_menu\_\_request\_scope\_check\_\_delta\_\_refresh\_notice\_\_selection\_\_response\_menu\_\_packet\_membership\_check}\\
$\to$ \texttt{worked\_example\_receipt\_interlock\_successor\_public\_request\_response\_menu\_\_request\_scope\_check\_\_delta\_\_refresh\_notice\_\_selection\_\_response\_menu\_\_packet\_membership\_check\_\_closure}.
\end{quote}

\paragraph{Worked stale-shortcut routing card.}
The paired-terminal shortcut is only useful if a later terse paper can also tell, immediately, which narrower or wider owner regains authority once that shortcut has gone stale.
\Cref{tab:workedexample-first-reopen-owner-card} pins that answer to the exact maintained worked objects rather than leaving it as free prose.
It is therefore the worked-example companion to the late-tail first-reopen-owner cards: same fail-closed habit, but grounded in one auditable route rather than only in the abstract owners.\cite{series:challengeanswerreviewmenus,series:publicrequestresponsepackets,series:publicrequestresponsepacketcarryforward,series:publicrequestresponsepacketdeltas,series:publicrequestresponsepacketrefreshnotices,series:publicrequestresponsepacketrefreshnoticenormalforms,series:publicrequestresponsepacketrefreshnoticeselections,series:publicrequestresponsepacketrefreshresponsemenus,series:publicrequestresponsepacketrefreshresponsepackets,series:publicrequestresponsepacketrefreshresponsepacketclosureverdicts,series:pairedterminalcitationdiscipline,series:pairedterminalcutoverrules}

\begin{table}[t]
\centering
\small
\begin{tabularx}{\linewidth}{@{}Y Y Y@{}}
\toprule
If the shortcut has gone stale because \ldots & First honest owner to reopen & Why this is first \\
\midrule
The live issue is exact clause location or a smaller semantic-piece support challenge & Citation docket / carrier slot \texttt{auditor\_compact\_diff\_summary\_answer\_carrier\_slot} together with the answer audit trail \texttt{auditor\_compact\_diff\_summary\_answer\_audit\_trail} & The answer-side terminal form has stopped one layer too early; the worked answer-side middle owners are live again before any source-only widening is honest. \\
Reveal order is the live issue & Answer escalation ladder \texttt{auditor\_compact\_diff\_summary\_answer\_escalation\_ladder} & This is the first owner that says how a referee would be shown more, so neither terminal shortcut may suppress it. \\
The first bounded omitted-support slice, its carried support surface, or its carryforward/delta basis is live & \texttt{worked\_example\_receipt\_interlock\_successor\_public\_request\_response\_menu\_\_request\_scope\_check} plus its \texttt{\_\_carryforward} / \texttt{\_\_delta} companions & The question has crossed the answer-side archive cut but has not yet become a visible-refresh family or successor-packet closure question. \\
The visible successor sentence is live only as emitted wrapper, canonical family, or selector choice & Refresh notice / normal form / selection beginning at \texttt{worked\_example\_receipt\_interlock\_successor\_public\_request\_response\_menu\_\_request\_scope\_check\_\_delta\_\_refresh\_notice} & The branch is still local to the visible wrapper layer, so the packet/closure capstone would already be too wide. \\
The live issue is the successor-facing follow-up class, exact packet, or branch closure verdict & Refresh response menu / packet / closure beginning at \texttt{worked\_example\_receipt\_interlock\_successor\_public\_request\_response\_menu\_\_request\_scope\_check\_\_delta\_\_refresh\_notice\_\_selection\_\_response\_menu} & This is the first request-side owner bucket that certifies the concrete successor-facing handoff, packet basis, and closed-versus-reopen posture. \\
The only live issue is whether the answer-side and request-side terminals may still travel together, or whether a later terse paper must stop, widen, or fail closed & Synthesis~74 for paired adjacency/reuse; Synthesis~75 for the stop-versus-widen judgment & Once the narrower answer-side/request-side owners are already fixed, the remaining live issue is the paired discipline or the cutover rule itself rather than any narrower maintained artifact. \\
\bottomrule
\end{tabularx}
\caption{Worked stale-shortcut routing card. This is the smallest worked map from one stale compact citation packet to the first narrower or wider owner that regains authority.}
\label{tab:workedexample-first-reopen-owner-card}
\end{table}

\paragraph{One tiny stale-shortcut drill.}
Suppose a later terse paper starts from the paired shortcut and then the referee asks \emph{``where exactly is that compact-diff sentence carried?''}
\Cref{tab:workedexample-first-reopen-owner-card} routes the note back to \texttt{auditor\_compact\_diff\_summary\_answer\_carrier\_slot} and \texttt{auditor\_compact\_diff\_summary\_answer\_audit\_trail}, not to the source-only branch.
If the next question instead becomes \emph{``which visible refresh sentence now governs the request-scope packet?''}, the same card reopens the refresh-notice / selector layer rather than jumping straight to the successor-facing packet capstone.
If the branch stays request-side but \nolinkurl{packet\_reuse\_reopens} fires, the same card pushes the paper down to the exact request-side carryforward / delta owner before any terminal shortcut is quoted again.
And if all narrower owners still agree but the only remaining live issue is \emph{``may I still cite both terminals together, or must I stop / widen / fail closed differently?''}, the same card widens to Synthesis~74--75 rather than pretending that the worked route alone owns that judgment.

\paragraph{Front-edge source-only contract/status card.}
Before any visible wrapper is emitted, later terse papers may need the exact worked front edge that licenses that wrapper at all.
\Cref{tab:workedexample-contract-status-card} therefore keeps the source-only request basis on one referee-facing table: the paper-facing request promise, the imported evidence-class vocabulary, the family-scoped completion basis, the within-release paper-level status token, and the successor carryforward verdict.
That keeps the Synthesis~56--60 split concrete in one worked surface: the contract owns what the paper promised, the evidence classes own the reusable omitted-support family names, the fulfillment certificates own whether each promised family has actually been delivered, the status envelope owns the paper-level token emitted from those family verdicts, and the carryforward envelope owns whether that token survives the concrete release pair unchanged, advances, or reopens.\cite{series:publicrequestcontracts,series:requestableevidenceclasses,series:requestfulfillmentcertificates,series:publicrequeststatusenvelopes,series:publicrequestcarryforward}

\begin{table}[t]
\centering
\footnotesize
\begin{tabularx}{\linewidth}{@{}p{0.24\linewidth}X@{}}
\toprule
Layer & Worked front-edge owner \\
\midrule
Paper-facing request promise & \nolinkurl{worked_example_receipt_interlock_public_request_contract} with request phrase ``supporting artifacts remain available on request,'' covering the public-anchor families \nolinkurl{public_status_sentence_family} and \nolinkurl{compact_diff_summary_family}. \\
Imported omitted-support family vocabulary & \nolinkurl{example_requestable_evidence_classes.json} with the shared class \nolinkurl{bundle_identity_and_validator_companions} and the family-local classes \nolinkurl{public_status_lineage_and_closure_support} and \nolinkurl{compact_diff_compare_and_replay_support}. \\
Family-scoped completion basis & \nolinkurl{public_status_sentence_request_fulfillment_certificate} and \nolinkurl{compact_diff_summary_request_fulfillment_certificate}, each with \nolinkurl{completion_verdict=complete}. \\
Within-release paper-level status token & \nolinkurl{worked_example_receipt_interlock_public_request_status_envelope} with \nolinkurl{public_request_status_token=complete}. \\
Successor carryforward verdict & \nolinkurl{worked_example_receipt_interlock_public_request_carryforward_envelope} with \nolinkurl{predecessor_public_request_status_token=complete}, \nolinkurl{successor_public_request_status_token=complete}, and \nolinkurl{carryforward_token=preserved_complete}. \\
Smallest honest front-edge stop & Promise identity $\to$ contract; omitted-support family meaning $\to$ evidence class; family-level delivered-versus-missing basis $\to$ fulfillment certificate; current paper-level token $\to$ status envelope; release-pair persistence/reopen judgment $\to$ carryforward envelope. \\
\bottomrule
\end{tabularx}
\caption{The worked source-only contract/status front edge. This is the smallest referee-facing card for recovering which promise, evidence-class vocabulary, family completion basis, paper-level status token, and release-pair carryforward verdict license the later visible wrapper layer.}
\label{tab:workedexample-contract-status-card}
\end{table}

\paragraph{One tiny contract-status drill.}
Suppose a later terse paper asks why the worked source-only promise may honestly compress to the paper-level token \nolinkurl{complete} before quoting any outward-facing sentence.
The first honest stop is not the visible notice layer at all but the status envelope \nolinkurl{worked_example_receipt_interlock_public_request_status_envelope}, because that is where the two family-scoped completion verdicts are compressed to one paper-level token.
If the live issue becomes why one of those family verdicts is complete rather than incomplete, widen exactly once more to the relevant fulfillment certificate; and if the live issue becomes what family meaning that certificate imported, widen exactly once more to the evidence-class ledger rather than stretching the certificate into a family-vocabulary note.
Only when the live issue becomes whether that paper-level token survives the concrete release pair unchanged should the route widen to \nolinkurl{worked_example_receipt_interlock_public_request_carryforward_envelope}.\cite{series:requestfulfillmentcertificates,series:publicrequeststatusenvelopes,series:publicrequestcarryforward}

\paragraph{Visible source-only notice-ladder card.}
Later terse papers may also need the concrete visible-wrapper ladder rather than the request-side terminal capstone.
\Cref{tab:workedexample-notice-ladder-card} therefore places the worked within-release and successor-facing source-only ladders side by side.
Each ladder keeps five adjacent owners explicit: the basis envelope and imported tokens, the selector and decision label, the chosen normal form, the emitted visible notice, and the first notice-keyed response menu that would answer a concrete follow-up.
That keeps the Synthesis~59--64 split concrete in one worked surface: the envelope emits the basis token, the selector chooses the family, the normal form owns the reusable sentence family, the notice owns the exact visible sentence, and the response menu owns the next follow-up handoff.\cite{series:publicrequeststatusenvelopes,series:publicrequestcarryforward,series:publicrequestnotices,series:publicrequestnoticenormalforms,series:publicrequestnoticeselections,series:publicrequestresponsemenus}

\begin{table}[t]
\centering
\footnotesize
\begin{tabularx}{\linewidth}{@{}p{0.21\linewidth}XX@{}}
\toprule
Layer & Within-release worked ladder & Successor-facing worked ladder \\
\midrule
Basis owner and imported tokens & \nolinkurl{worked_example_receipt_interlock_public_request_status_envelope} with \nolinkurl{public_request_status_token=complete}. & \nolinkurl{worked_example_receipt_interlock_public_request_carryforward_envelope} with \nolinkurl{public_request_status_token=complete} and \nolinkurl{carryforward_token=preserved_complete}. \\
Selector and decision & \nolinkurl{worked_example_receipt_interlock_within_release_public_request_notice_selection} with \nolinkurl{select_within_release_complete_notice_normal_form}. & \nolinkurl{worked_example_receipt_interlock_successor_public_request_notice_selection} with \nolinkurl{select_successor_preserved_complete_notice_normal_form}. \\
Chosen normal form & \nolinkurl{worked_example_receipt_interlock_within_release_complete_public_request_notice_normal_form}. & \nolinkurl{worked_example_receipt_interlock_successor_preserved_complete_public_request_notice_normal_form}. \\
Emitted notice / sentence & \nolinkurl{worked_example_receipt_interlock_public_request_status_notice}: ``For this source-only release, the worked note's paper-level public request promise is complete and supporting artifacts remain available on request.'' & \nolinkurl{worked_example_receipt_interlock_public_request_carryforward_notice}: ``Across the canned successor release, that complete paper-level public request status carries forward unchanged.'' \\
First concrete follow-up handoff & \nolinkurl{worked_example_receipt_interlock_within_release_public_request_response_menu}. & \nolinkurl{worked_example_receipt_interlock_successor_public_request_response_menu}. \\
Smallest honest wrapper stop & Sentence identity $\to$ emitted notice; reusable family $\to$ normal form; family-choice why $\to$ selector. & Sentence identity $\to$ emitted notice; reusable family $\to$ normal form; carryforward-family choice why $\to$ selector. \\
\bottomrule
\end{tabularx}
\caption{The worked visible source-only notice ladders, side by side. This is the smallest referee-facing card for recovering which basis tokens, selectors, normal forms, emitted notices, and first follow-up menus govern the within-release and successor-facing wrappers without reopening the whole source-only routing tail.}
\label{tab:workedexample-notice-ladder-card}
\end{table}

\paragraph{One tiny visible-wrapper drill.}
Suppose a later terse paper asks why the worked successor sentence may say ``carries forward unchanged.''
The first honest stop is not the emitted notice alone but the successor selector \nolinkurl{worked_example_receipt_interlock_successor_public_request_notice_selection}, because that selector is where the imported token bundle \nolinkurl{public_request_status_token=complete} plus \nolinkurl{carryforward_token=preserved_complete} is forced to the family \nolinkurl{worked_example_receipt_interlock_successor_preserved_complete_public_request_notice_normal_form}.
By contrast, if the only question is what visible sentence readers actually saw, the honest stop is one layer later at \nolinkurl{worked_example_receipt_interlock_public_request_carryforward_notice}; and if the question becomes ``which concrete follow-up owner should answer next?'', the route widens exactly once more to \nolinkurl{worked_example_receipt_interlock_successor_public_request_response_menu}.\cite{series:publicrequestnotices,series:publicrequestnoticenormalforms,series:publicrequestnoticeselections,series:publicrequestresponsemenus}


\paragraph{Visible successor-refresh ladder card.}
Later terse papers may also need the concrete successor-refresh wrapper ladder rather than only the request-side terminal capstone.
\Cref{tab:workedexample-refresh-ladder-card} therefore places the exact worked refresh chain on one referee-facing table.
That keeps the Synthesis~67--73 split concrete in one place: the delta ledger owns which field families refreshed, the selector chooses the canonical visible reissue family, the normal form owns the reusable wrapper family, the emitted refresh notice owns the exact visible sentence, the successor-facing response menu owns the next follow-up handoff, the successor-facing packet owns the exact bounded slice, and the closure verdict owns whether that branch is currently a stable finite public basis.\cite{series:publicrequestresponsepacketdeltas,series:publicrequestresponsepacketrefreshnotices,series:publicrequestresponsepacketrefreshnoticenormalforms,series:publicrequestresponsepacketrefreshnoticeselections,series:publicrequestresponsepacketrefreshresponsemenus,series:publicrequestresponsepacketrefreshresponsepackets,series:publicrequestresponsepacketrefreshresponsepacketclosureverdicts}

\begin{table}[t]
\centering
\footnotesize
\begin{tabularx}{\linewidth}{@{}p{0.24\linewidth}X@{}}
\toprule
Layer & Successor-facing worked ladder \\
\midrule
Basis owner and imported reuse tokens & \nolinkurl{worked_example_receipt_interlock_successor_public_request_response_menu__request_scope_check__carryforward} together with \nolinkurl{worked_example_receipt_interlock_successor_public_request_response_menu__request_scope_check__delta}, fixing \nolinkurl{packet_reuse_status=refresh_in_place} and \nolinkurl{delta_token=refresh_fields}. \\
Selector and decision & \nolinkurl{worked_example_receipt_interlock_successor_public_request_response_menu__request_scope_check__delta__refresh_notice__selection} with \nolinkurl{select_successor_carryforward_request_scope_check_refresh_fields_packet_refresh_notice_normal_form}. \\
Chosen normal form & \nolinkurl{worked_example_receipt_interlock_successor_public_request_response_menu__request_scope_check__delta__refresh_notice__normal_form}. \\
Emitted refresh notice / sentence & \nolinkurl{worked_example_receipt_interlock_successor_public_request_response_menu__request_scope_check__delta__refresh_notice}: ``Under successor release worked-example-draft-194b, the request-scope check response packet for this successor-carryforward source-only notice remains the same bounded handoff slice; only the basis companions, selected entry excerpts, and digest bindings families refresh in place.'' \\
First successor-facing follow-up handoff & \nolinkurl{worked_example_receipt_interlock_successor_public_request_response_menu__request_scope_check__delta__refresh_notice__selection__response_menu}. \\
Exact successor-facing packet & \nolinkurl{worked_example_receipt_interlock_successor_public_request_response_menu__request_scope_check__delta__refresh_notice__selection__response_menu__packet_membership_check}. \\
Closure verdict and reopen floor & \nolinkurl{worked_example_receipt_interlock_successor_public_request_response_menu__request_scope_check__delta__refresh_notice__selection__response_menu__packet_membership_check__closure} with \nolinkurl{closure_status=closed_finite_basis}, \nolinkurl{closure_basis_scope=public_paths_only}, and reopen triggers \nolinkurl{followup_taxonomy_expands}, \nolinkurl{packet_reuse_reopens}, and \nolinkurl{visible_refresh_sentence_rebinds_to_new_delta_family}. \\
Smallest honest refresh stop & Sentence identity $\to$ emitted refresh notice; reusable wrapper family $\to$ normal form; family-choice why $\to$ selector; next follow-up owner $\to$ response menu; exact successor-facing public basis or current branch closure $\to$ packet plus closure verdict. \\
\bottomrule
\end{tabularx}
\caption{The worked visible successor-refresh ladder for the request-scope branch. This is the smallest referee-facing card for recovering which delta basis, selector, normal form, emitted refresh notice, successor-facing handoff, exact packet, and closure verdict govern the current visible-refresh branch without reopening the whole successor tail.}
\label{tab:workedexample-refresh-ladder-card}
\end{table}

\paragraph{One tiny visible-refresh drill.}
Suppose a later terse paper asks why the worked successor branch may still answer \emph{packet membership} without reopening the request-scope slice.
If the live issue is only the sentence readers actually saw, the honest stop is \nolinkurl{worked_example_receipt_interlock_successor_public_request_response_menu__request_scope_check__delta__refresh_notice}.
If the live issue is instead why that visible sentence is the canonical one under the shipped delta token, the honest stop is one layer earlier at the selector \nolinkurl{worked_example_receipt_interlock_successor_public_request_response_menu__request_scope_check__delta__refresh_notice__selection}.
If the live issue becomes the exact successor-facing packet that should travel next, widen exactly once more to \nolinkurl{worked_example_receipt_interlock_successor_public_request_response_menu__request_scope_check__delta__refresh_notice__selection__response_menu__packet_membership_check}; and if the live issue becomes whether that branch already closes there, widen exactly once more to the closure verdict beside it.
As soon as any named reopen trigger fires, the refresh-ladder shortcut goes stale and the paper must fail closed to the named reopening owner rather than keep quoting the packet/closure pair as if the branch were unchanged.\cite{series:publicrequestresponsepacketrefreshnotices,series:publicrequestresponsepacketrefreshnoticeselections,series:publicrequestresponsepacketrefreshresponsemenus,series:publicrequestresponsepacketrefreshresponsepackets,series:publicrequestresponsepacketrefreshresponsepacketclosureverdicts}

The answer-audit trail keeps the semantic pieces just as small.
For the shipped auditor compact-diff route, only three piece-level terminals are needed:
\begin{center}
\small
\begin{tabularx}{0.97\linewidth}{@{}Y Y Y@{}}
\toprule
Piece & Terminal owner & Why stop there \\
\midrule
diff\_delta\_piece & \texttt{example\_compare\_report.json} : \texttt{observed\_diffs} & This is the smallest numeric-delta owner. \\
recompute\_piece & \texttt{example\_verifier\_report.json} : \texttt{recomputed} & This is the smallest replay-facing recomputation witness. \\
changed\_family\_piece & \texttt{example\_successor\_continuity\_verdict.json} : \texttt{changed\_line\_items} & This is the smallest changed-family owner used by downstream maintenance notes. \\
\bottomrule
\end{tabularx}
\end{center}
\vspace{-0.5em}

\noindent The important discipline is that nothing widens just because a larger artifact exists.
The walk stops at the earliest sufficient owner for the actual question being asked.
That is the operational payoff of the whole spine: the example is auditable precisely because the drill is narrow, explicit, and owner-stable.

\section{Illustrative arithmetic slice, not a channel certificate}
\label{sec:numerical}
This note retains one bundle-backed \emph{arithmetic} slice so the public prose and the maintained validator use the same concrete inputs. The receipt itself labels the slice
\nolinkurl{semantic_scope=illustrative_tiered_observation_attenuation_not_mceq},
\nolinkurl{composition_evidence_status=assumption_only_no_joint_channel_witness}, and
\nolinkurl{publication_eligible=false}.
Those fields are part of the result, not administrative decoration. The shipped verifier proves that the stated formulas were evaluated correctly; it does not prove that scalar activation probabilities describe the deployment's observation channel.\cite{series:auditorreplay,series:obsatten}

\paragraph{Attenuation under an explicit common-mixture wrapper.}
The formula
\[
b_{\mathrm{att}}=\log_2\!\big((1-p)+p2^b\big)
\]
is valid when the observed channel is a source-bound common-mixture/erasure wrapper: the activation gate is independent of the secret and the inner transcript, the inactive output has a secret-independent law, and the conditional active-channel leakage is at most $b$. Equal per-secret activation rates alone do not suffice.
The maintained bundle stores $p_{\mathrm{obs}}\le0.02$ as an assumption-only sensitivity input rather than a joint-channel witness. Its tier values therefore remain arithmetic regression values. Likewise, the fallback-rate mixture and the selection-tax arithmetic below are non-authorizing unless their path-selection and joint-channel premises are separately witnessed.

\begin{table}[t]
\centering
\small
\begin{tabularx}{0.96\linewidth}{@{}Y Y Y@{}}
\toprule
Maintained quantity & Arithmetic inputs and missing premise & Recomputed value (bits)\\
\midrule
\texttt{fallback\_contact\_surface} & $b=1.40$, $p=0.02$; requires a source-bound common-mixture channel & 0.04653\\
\texttt{fallback\_tiered\_vector} & CSET/TIME/CONG values $(1.40,0.30,0.15)$ at $p=0.02$; no joint-tier witness & 0.05635\\
\texttt{path\_selection\_indicator} & arithmetic tax $\log_2 R_{\pi}$ with $R_{\pi}\le1.2$; requires a secret/path-selector channel witness & 0.26303\\
Illustrative total arithmetic & $\widehat p_{\mathrm{fb}}=0.05$, $b_{\mathrm{primary}}=0$, $b_{\mathrm{fallback}}=0.05635$, plus selection tax & 0.26590\\
\bottomrule
\end{tabularx}
\caption{A non-authorizing arithmetic regression from the maintained bundle. Correct recomputation does not establish the missing joint observation-channel premises.}
\label{tab:workedexample-bundle-backed-slice}
\end{table}

\paragraph{Why the slice remains useful.}
The point is to show one exact replay chain from declared arithmetic inputs to one output and then stop. It catches evaluator drift, rounding changes, and broken artifact bindings. It cannot turn an assumption-only scalar into a deployment privacy certificate. A publication-eligible successor must supply source-bound channel evidence for the activation gate, the tier joint law, the path selector, and the stopping/repetition semantics; until then, later papers may quote these values only as illustrative arithmetic and must preserve the three fail-closed receipt fields above.\cite{series:fallbacktax,series:traceifaces,series:stataudit,series:anytimespend}

\section{Optional profiling variants (three one-paragraph receipts)}
\paragraph{Prefix-fetch equalization projection.}
A variant receipt replaces the fallback ENF with a prefix-fetch projection (coarser query privacy via requesting a bucket/prefix rather than a single key).
The interface change is surfaced \emph{only} by a different ENF digest and different evidence digest; the rest of the receipt machinery is unchanged.\cite{series:prefixcapacities,series:stataudit}
Concretely, \nolinkurl{example\_receipt\_prefixfetch\_variant.json} adds \nolinkurl{profiling\_equalization\_prefixfetch} with \nolinkurl{evidence\_id=sha256:887e6706ae2e1b982f28dbafdbe7796e59b06ab2dbeb5f4db43fe9e667268f61}, \nolinkurl{budget\_value=(eta\_bits=0.292597, delta=0.004483)}, replay hook \nolinkurl{eval-profiling-prefixfetch-v1}, and support bundle \nolinkurl{bundle://worked-example/profiling-prefixfetch}.
The state-conditioned variant \nolinkurl{example\_receipt\_prefixfetch\_statecond\_variant.json} instead adds \nolinkurl{profiling\_equalization\_prefixfetch\_statecond} with \nolinkurl{evidence\_id=sha256:348176a3e0bdaede5a67f8497ed1c31effd1d64f34eae12e9e9b46154b0b6bd3}, \nolinkurl{budget\_value=(eta\_bits=0.459923, delta=0.007146)}, replay hook \nolinkurl{eval-profiling-prefixfetch-statecond-v1}, and support bundle \nolinkurl{bundle://worked-example/profiling-prefixfetch-statecond}.

\paragraph{Coarsening a large-alphabet projection.}
A second variant receipt demonstrates safe coarsening of a large-alphabet contact projection into an audit-feasible partition.
Again, the only user-visible change is via the ENF and evidence digests, plus the resulting budget number.\cite{series:coarsening}
The maintained coarsening variant \nolinkurl{example\_receipt\_contact\_coarsened\_variant.json} adds \nolinkurl{profiling\_equalization\_contact\_coarsened} with \nolinkurl{evidence\_id=sha256:22666e9853213a11356f4b863aa9e0dba88bcdd56306ac345acd80e1f09d36fa}, \nolinkurl{coarsening\_map\_id=sha256:4c447e40d9e29e41ffa313f41edd5b2d54ad296c5a9f58b01d26ac220db5bbc1}, \nolinkurl{budget\_value=(eta\_bits=0.391684, delta=0.0)}, replay hook \nolinkurl{eval-profiling-contact-coarsened-v1}, and support bundle \nolinkurl{bundle://worked-example/profiling-contact-coarsened}.
If the partition map is selected from data (rather than fixed ex ante), the selection discipline note applies: treat the map family as witness-indexed or separate pilot selection from frozen certification under change control.\cite{series:selectiondiscipline}

The base owner-map row is explicit before the variants: \nolinkurl{profiling_equalization} in \nolinkurl{example_line_item_owner_map.json} now carries \nolinkurl{materialization_status=base_receipt_current_cut}, \nolinkurl{receipt_file=example_receipt.json}, \nolinkurl{support_bundle_id=bundle://worked-example/profiling-eq}, and public objects \nolinkurl{example_receipt.json}, \nolinkurl{example_profiling_evidence.json}, \nolinkurl{example_replay_plans.json}, and \nolinkurl{example_support_bundle_map.json}. This makes the maintained base profiling card symmetric with its replay/support bundle instead of leaving the owner-map row thinner than the receipt claim.

The owner-map side is explicit for the variants as well: the three \nolinkurl{variant_entries} in \nolinkurl{example_line_item_owner_map.json} carry \nolinkurl{materialization_status=optional_variant_receipt_current_cut}, \nolinkurl{receipt_variant_file}, \nolinkurl{evidence_id}, \nolinkurl{replay_plan_id}, \nolinkurl{plan_spec_id}, \nolinkurl{support_bundle_id}, \nolinkurl{support_consumed_by}, \nolinkurl{support_required_artifacts}, \nolinkurl{support_pointers}, and \nolinkurl{public_objects}. The prefix row names \nolinkurl{example_receipt_prefixfetch_variant.json}; the state-conditioned row names \nolinkurl{example_receipt_prefixfetch_statecond_variant.json}; and the contact-coarsened row names \nolinkurl{example_receipt_contact_coarsened_variant.json}. The verifier side now mirrors those optional rows under \nolinkurl{optional_variant_plans_checked}, where each optional row is separate from base \nolinkurl{plans_checked} and repeats the variant receipt file, evidence id, replay plan spec, support bundle, support-required artifacts, support pointers, public objects, and \nolinkurl{public_object_digest_entries}. The series-spine side mirrors those rows under \nolinkurl{optional_variant_spines}, with each row carrying \nolinkurl{variant_sentence_normal_form} plus receipt, owner-map, replay, and support anchors, so these optional receipt variants are mechanically quotable without being folded into the base \nolinkurl{line_item_spines} ladder. This keeps optional profiling-variant receipts visibly separated from base-receipt materialization, base verifier \nolinkurl{plans_checked}, and owner-map-only route labels.

\section{Takeaways and open edges (short)}
\begin{itemize}[leftmargin=*]
\item A referee-facing object can stay tiny: digests + a handful of line items + replay hooks.
\item ``Same name, different semantics'' is handled by digest IDs (ENF/TW/state), not by prose.
\item Notation (keys/state/transcripts, projections) is standardized in Synthesis~8; this note only instantiates it.\cite{series:notation}
\item Guarded endpoints are essential: a primary-path claim that depends on non-collusion must ship its $\rho$.
\item Tier vectors and selection tax compile to a single per-lookup MaxL (or odds) parameter that downstream tools can consume.\cite{series:endpointbridge}
\item Remaining research: robust horizon accounting under stateful dependence and operator-safe drift discipline; the state series and accounting rulebook are the intended import points.\cite{series:state1,series:accountingrulebook}
\end{itemize}

\begin{thebibliography}{99}\small


\bibitem{rfc8785}
A.~Rundgren.
\newblock JSON Canonicalization Scheme (JCS).
\newblock RFC~8785, IETF, 2020.

\bibitem{series:artifactpolicy}
Anonymous.

\newblock Artifact and Disclosure Policy for the One-True-Archive.

\newblock Synthesis~6, The-One-True-Archive (2026).

\bibitem{series:downstreamnormalforms}
Anonymous.
\newblock Downstream Normal Forms for Successor Maintenance: Summary, Support, Citation, Quote, and Clause Discipline.
\newblock Series synthesis note (Synthesis~31), 2026. (This archive.)

\bibitem{series:mathcrosswalk}
Anonymous.

\newblock Mathematics Backbone Crosswalk for Anonymous-DHT Receipts: What the five published wiki notes export, and where to cite them.

\newblock Series synthesis note (Synthesis~24), 2026. (This archive.)


\bibitem{series:schedretries}
Anonymous.

\newblock Schedule-Only Retries in Anonymous DHTs: When Retry Logic Becomes a Linkability Channel.

\newblock Operational series Paper~A, The-One-True-Archive (2026).

\bibitem{series:operationalB}
Anonymous.

\newblock Auditable Trace Privacy for Anonymous DHTs: Plan Pinning, PTLs, and Epoch Conformance Receipts.

\newblock Operational series Paper~B, The-One-True-Archive (2026).

\bibitem{series:advantagecontracts}
Anonymous.

\newblock Advantage Contracts and Stealth Audits for Anonymous DHTs.

\newblock Evaluation series Paper~2, The-One-True-Archive (2026).
\bibitem{series:seriesspines}
Anonymous.

\newblock Series Spines from Mathematical Roots to Referee Handoffs: A Non-Redundant Bridge from Backbones, Receipts, Replay, Release Evolution, and Review Menus.

\newblock Series synthesis note (Synthesis~55), 2026. (This archive.)

\bibitem{series:deployedsurfaces}
Anonymous.

\newblock Deployed Observation Surfaces for Anonymous DHTs: Control-plane knobs in IPFS/libp2p delegated routing and OHTTP evolution.

\newblock Series synthesis note (Synthesis~30), 2026. (This archive.)

\bibitem{series:receiptitems}
Anonymous.

\newblock Receipt Line Items for Anonymous-DHT Privacy Claims: Minimal Tuple Schema and Drift Semantics.

\newblock Series synthesis note (Synthesis~12), 2026. (This archive.)

\bibitem{series:notation}
Anonymous.

\newblock Notation and a Minimal Model for the Series.

\newblock Series synthesis note (Synthesis~8), 2026. (This archive.)

\bibitem{series:traceifaces}
Anonymous.

\newblock Trace Interfaces for Anonymous DHTs.

\newblock Synthesis~3, The-One-True-Archive (2026).

\bibitem{series:threatwindows}
Anonymous.

\newblock Threat Models and Repetition Windows for Anonymous DHTs.

\newblock Synthesis~4, The-One-True-Archive (2026).

\bibitem{series:endpointbridge}
Anonymous.

\newblock Endpoint Metrics Bridge for Anonymous Overlays: Posterior-Mass Inflation, PML, MaxL, and True Odds.

\newblock Synthesis~5, The-One-True-Archive (2026).

\bibitem{series:oddsinflation}
Anonymous.

\newblock Posterior-Mass Inflation Anonymity: Pointwise Maximal Leakage, Composition, and the True-Odds Boundary.

\newblock Anonymity~A, The-One-True-Archive (2026).

\bibitem{series:profiling}
Anonymous.

\newblock Anonymous DHT Profiling and the Equalization Contract.

\newblock Anonymity~B, The-One-True-Archive (2026).

\bibitem{series:primarysepmanifest}
Anonymous.

\newblock Primary-Path Separation Manifests: Guard Objects for Non-Collusion Assumptions in Reader-Privacy Receipts.

\newblock Synthesis~22, The-One-True-Archive (2026).

\bibitem{series:probsep}
Anonymous.

\newblock Probabilistic Separation Budgets: Separation Failure as a Receipt-Grade $(\varepsilon,\delta)$ Adjunct.

\newblock Synthesis~23, The-One-True-Archive (2026).

\bibitem{series:obsatten}
Anonymous.

\newblock Observation Attenuation for Repeated-Use Privacy Budgets.

\newblock Synthesis~11, The-One-True-Archive (2026).


\bibitem{series:planstateequiv}
Anonymous.

\newblock Digest-Bound Replay Plans and State Declarations: Equivalence, Minimality, and Change Control.

\newblock Synthesis~18, The-One-True-Archive (2026).

\bibitem{series:userminiface}
Anonymous.

\newblock User-Verifiable Receipt Interfaces.

\newblock Synthesis~13, The-One-True-Archive (2026).

\bibitem{series:fallbacktax}
Anonymous.

\newblock Adaptive Fallback as a Witness Stream: The Path-Selection Tax.

\newblock Synthesis~14, The-One-True-Archive (2026).

\bibitem{series:tieredobs}
Anonymous.

\newblock Tiered Observation Vectors for Anonymous-DHT Receipts.

\newblock Synthesis~15, The-One-True-Archive (2026).

\bibitem{series:abomoinl}
Anonymous.

\newblock ABOM and OINL for Anonymity Claims.

\newblock Anonymity series Paper~C, The-One-True-Archive (2026).

\bibitem{series:bossfightB}
Anonymous.

\newblock AnonDHT in the Boss Fight Regime: Compiling a MaxL Budget into Lookup Knobs.

\newblock Boss Fight series Paper~B, The-One-True-Archive (2026).


\bibitem{series:bossfightA}
Anonymous.

\newblock The Boss Fight of Dummy Cover for Anonymous-DHT Lookups.

\newblock Boss Fight series Paper~A, The-One-True-Archive (2026).

\bibitem{series:bossfightC}
Anonymous.

\newblock Proof-Carrying Budgets for Anonymous DHTs.

\newblock Boss Fight series Paper~C, The-One-True-Archive (2026).

\bibitem{series:evalnotary}
Anonymous.

\newblock Notarized Anonymity Certificates Under Retries.

\newblock Evaluation series Paper~1, The-One-True-Archive (2026).

\bibitem{series:releaseproperty}
Anonymous.

\newblock Release Properties, Lineage Tags, and Ledger Receipts for Anonymous Deployments.

\newblock Release and destination series Paper~A, The-One-True-Archive (2026).


\bibitem{series:prefixcapacities}
Anonymous.

\newblock Prefix Capacities and Separation Cuts for Shared-Kernel Termination-Time Hiding.

\newblock Published mathematics note, 2026-01-25. (This archive: \texttt{published/2026-01-25\_prefix\_capacities}.)

\bibitem{series:stoptimenf}
Anonymous.

\newblock Stop-Time Padding Addendum, Max-Mixture Separations and Tail-Sign Deadline Normal Forms.

\newblock Published mathematics addendum, 2026-01-26. (This archive: \texttt{published/2026-01-26\_stop\_time\_padding\_addendum}.)

\bibitem{series:state1}
Anonymous.

\newblock State-Dependent Anonymity for Anonymous DHTs.

\newblock AnonDHT state series Paper~1, The-One-True-Archive (2026).

\bibitem{series:state1A}
Anonymous.

\newblock State-Dependent Anonymity (Addendum): Calibration Vignettes and Time-to-Identification Conversions.

\newblock AnonDHT state series Paper~1A, The-One-True-Archive (2026).

\bibitem{series:mucc}
Anonymous.

\newblock Marginally Uniform Committee Contact in Anonymous DHT Lookups: Availability Floors, Perfect-Leakage Counterexamples, and Abort-Safe Equalization.

\newblock AnonDHT state series Paper~2, The-One-True-Archive (2026).

\bibitem{series:state3}
Anonymous.

\newblock Closed-View Auditing for Stateful Anonymity: Odometers, Alarms, and Machine-Checkable CPPCs.

\newblock AnonDHT state series Paper~3, The-One-True-Archive (2026).

\bibitem{series:state3A}
Anonymous.

\newblock Closed-View Auditing and CPPCs (Addendum): Multi-Stream Monitoring Notes and Schema Excerpts.

\newblock AnonDHT state series Paper~3A, The-One-True-Archive (2026).

\bibitem{series:mceq}
Anonymous.

\newblock MC-EQ and CoverSketch: Destination-Privacy Equalization for Anonymous-DHT Lookups.

\newblock Release and destination series Paper~B, The-One-True-Archive (2026).

\bibitem{series:eval1A}
Anonymous.

\newblock Notarized Anonymity Certificates (Addendum): Witness Hardening and Optional Checkpoint Notes.

\newblock Evaluation series Paper~1A, The-One-True-Archive (2026).

\bibitem{series:eval2}
Anonymous.

\newblock Advantage Contracts and Stealth Audits for Anonymous DHT Evaluation.

\newblock Evaluation series Paper~2, The-One-True-Archive (2026).

\bibitem{series:eval3}
Anonymous.

\newblock Calibration Recipes for Anonymous DHTs.

\newblock Evaluation series Paper~3, The-One-True-Archive (2026).

\bibitem{series:condcert}
Anonymous.

\newblock Conditional Per-Step Budget Certificates: A Receipt-Grade Interface for Adaptive Horizon Claims.

\newblock Synthesis~21, The-One-True-Archive (2026).

\bibitem{series:accountingrulebook}
Anonymous.

\newblock Receipt Accounting and Horizon Composition for Anonymous-DHT Budgets.

\newblock Synthesis~19, The-One-True-Archive (2026).


\bibitem{series:stataudit}
Anonymous.

\newblock Discrete Equalization Audits from Logs: Menu-Level Finite-Sample TV Certificates.

\newblock Synthesis~25, The-One-True-Archive (2026).

\bibitem{series:anytimespend}
Anonymous.

\newblock Anytime Statistical Spending for Menu Audits: Change-Control-Safe Confidence Budgets.

\newblock Synthesis~26, The-One-True-Archive (2026).

\bibitem{series:coarsening}
Anonymous.

\newblock Coarsening and Large-Alphabet Projections for Log Audits: Safe Partitioning Rules that Keep Certificates Sample-Feasible.

\newblock Synthesis~28, The-One-True-Archive (2026).

\bibitem{series:selectiondiscipline}
Anonymous.

\newblock Selection Discipline for Data-Dependent Bucket Maps.

\newblock Synthesis~29, The-One-True-Archive (2026).

\bibitem{series:statusenvelopes}
Anonymous.

\newblock Successor Status Envelopes and Lineage Notices for Release Maintenance: Narrow Status Tokens, Reader Pointer Order, and Public Answers.

\newblock Series synthesis note (Synthesis~34), 2026. (This archive.)

\bibitem{series:releasespine}
Anonymous.

\newblock Stable Release Spine Contracts for Successor Maintenance: Stage Ownership, Handoff Rules, and Auditable Walk-Throughs.

\newblock Series synthesis note (Synthesis~32), 2026. (This archive.)

\bibitem{series:releaseclosure}
Anonymous.

\newblock Release Obligations and Closure Ledgers for Successor Maintenance: Publication Duties, Package Closure, and Narrow Verdict Boundaries.

\newblock Series synthesis note (Synthesis~33), 2026. (This archive.)

\bibitem{series:challengeroutes}
Anonymous.

\newblock Challenge Routes and Audience Profiles for Successor Maintenance: Audience-Keyed Start Objects and Escalation Order.

\newblock Series synthesis note (Synthesis~36), 2026. (This archive.)

\bibitem{series:challengestops}
Anonymous.
\newblock Challenge Stop Profiles for Successor Maintenance: Question Families, Semantic Pieces, and Approved Terminal Fields.
\newblock Series synthesis note (Synthesis~37), 2026. (This archive.)

\bibitem{series:auditorreplay}
Anonymous.

\newblock Auditor Replay for Anonymous-DHT Receipt Line Items: A Mechanical Checklist for Digest-Bound Claims.

\newblock Series synthesis note (Synthesis~20), 2026. (This archive.)


\bibitem{series:challengecoverage}
Anonymous.
\newblock Challenge Coverage Certificates for Successor Maintenance: Audience Coverage, Surface-Carried Pieces, and Support-Only Hooks.
\newblock Series synthesis note (Synthesis~39), 2026. (This archive.)

\bibitem{series:challengesurfacehooks}
Anonymous.
\newblock Challenge Surface Hooks for Successor Maintenance: Piece-to-Segment Realization and Shared Surface Bindings.
\newblock Series synthesis note (Synthesis~40), 2026. (This archive.)

\bibitem{series:challengeterminalwitnesses}
Anonymous.
\newblock Challenge Terminal Witness Ledgers for Successor Maintenance: Narrow Stop Ownership for Piece Hooks and Support-Only Escalation.
\newblock Series synthesis note (Synthesis~41), 2026. (This archive.)

\bibitem{series:challengecoveragewitnesscores}
Anonymous.
\newblock Challenge Coverage Witness Cores for Successor Maintenance: Audience-Invariant Certification Dependencies for One Coverage Claim.
\newblock Series synthesis note (Synthesis~44), 2026. (This archive.)

\bibitem{series:challengecoveragewitnesspacks}
Anonymous.
\newblock Challenge Coverage Witness Packs for Successor Maintenance: Aggregate Route-and-Branch Certification Dependencies for One Coverage Claim.
\newblock Series synthesis note (Synthesis~42), 2026. (This archive.)

\bibitem{series:challengecoveragewitnessslices}
Anonymous.
\newblock Challenge Coverage Witness Slices for Successor Maintenance: Audience-Minimal Route-and-Branch Certification Deltas for One Coverage Claim.
\newblock Series synthesis note (Synthesis~43), 2026. (This archive.)

\bibitem{series:challengeanswercapsules}
Anonymous.
\newblock Challenge Answer Capsules for Successor Maintenance: Audience-Question Minimal Answer Bundles from Certified Piece Hooks.
\newblock Series synthesis note (Synthesis~45), 2026. (This archive.)

\bibitem{series:challengeanswersheets}
Anonymous.
\newblock Challenge Answer Sheets for Successor Maintenance: Audience-Question Piece-Complete Handoffs from Answer Capsules.
\newblock Series synthesis note (Synthesis~46), 2026. (This archive.)

\bibitem{series:challengeanswercards}
Anonymous.
\newblock Challenge Answer Cards for Successor Maintenance: Audience-Question Released Answers from Answer Sheets and Sentence Locks.
\newblock Series synthesis note (Synthesis~47), 2026. (This archive.)

\bibitem{series:challengeanswercarrierslots}
Anonymous.
\newblock Challenge Answer Carrier Slots for Successor Maintenance: Exact Shipped Fields and Reusable Clauses from Answer Cards.
\newblock Series synthesis note (Synthesis~48), 2026. (This archive.)

\bibitem{series:challengeansweraudittrails}
Anonymous.
\newblock Challenge Answer Audit Trails for Successor Maintenance: Leftward Reader Walks from Carrier Slots to Narrow Terminals.
\newblock Series synthesis note (Synthesis~49), 2026. (This archive.)

\bibitem{series:challengeanswercitationdockets}
Anonymous.
\newblock Challenge Answer Citation Dockets for Successor Maintenance: Smallest-Sufficient Late-Stage Citation Choices from Capsules to Audit Trails.
\newblock Series synthesis note (Synthesis~50), 2026. (This archive.)

\bibitem{series:challengeanswerpublicationprofiles}
Anonymous.
\newblock Challenge Answer Publication Profiles for Successor Maintenance: Budgeted Reuse Modes from Answer Cards, Sentence Locks, and Citation Dockets.
\newblock Series synthesis note (Synthesis~51), 2026. (This archive.)

\bibitem{series:challengeanswerdisclosurepackets}
Anonymous.
\newblock Challenge Answer Disclosure Packets for Successor Maintenance: Public Pointers and Requestable Audit Slices under Space Pressure.
\newblock Series synthesis note (Synthesis~52), 2026. (This archive.)

\bibitem{series:challengeanswerescalationladders}
Anonymous.
\newblock Challenge Answer Escalation Ladders for Successor Maintenance: Ordered Brief-to-Audit Reveal Policies from Publication Profiles and Disclosure Packets.
\newblock Series synthesis note (Synthesis~53), 2026. (This archive.)

\bibitem{series:challengebranches}
Anonymous.
\newblock Challenge Branch Maps for Successor Maintenance: Piece-Keyed Escalation from Audience Routes to Approved Terminal Fields.
\newblock Series synthesis note (Synthesis~38), 2026. (This archive.)

\bibitem{series:challengeanswerreviewmenus}
Anonymous.
\newblock Challenge Answer Review Menus for Successor Maintenance: Smallest-Sufficient Handoffs for Common Referee Follow-Up Classes.
\newblock Series synthesis note (Synthesis~54), 2026. (This archive.)

\bibitem{series:publicrequestcontracts}
Anonymous.
\newblock Public Request Contracts for Source-Only Releases: Exact Paper-Level Handoffs from Boilerplate Policy to Maintained Packet Cuts.
\newblock Series synthesis note (Synthesis~56), 2026. (This archive.)

\bibitem{series:requestableevidenceclasses}
Anonymous.
\newblock Requestable Evidence Classes for Source-Only Releases: Stable Names for Omitted-Support Families behind Packet Cuts and Public Request Contracts.
\newblock Series synthesis note (Synthesis~57), 2026. (This archive.)


\bibitem{series:publicrequeststatusenvelopes}
Anonymous.
\newblock Public Request Status Envelopes for Source-Only Releases: Paper-Level Tokens Compressing Family Completion without Re-owning the Promise.
\newblock Series synthesis note (Synthesis~59), 2026. (This archive.)

\bibitem{series:publicrequestcarryforward}
Anonymous.
\newblock Public Request Carryforward for Source-Only Successor Releases: When Paper-Level Request Status Persists, Advances, or Reopens across Release Evolution.
\newblock Series synthesis note (Synthesis~60), 2026. (This archive.)

\bibitem{series:publicrequestnotices}
Anonymous.
\newblock Public Request Notices for Source-Only Releases: Outward-Facing Sentences for Within-Release Status and Cross-Release Carryforward.
\newblock Series synthesis note (Synthesis~61), 2026. (This archive.)

\bibitem{series:publicrequestnoticenormalforms}
Anonymous.
\newblock Public Request Notice Normal Forms for Source-Only Releases: Canonical Sentence Families for Within-Release Status and Cross-Release Carryforward Notices.
\newblock Series synthesis note (Synthesis~62), 2026. (This archive.)

\bibitem{series:publicrequestnoticeselections}
Anonymous.
\newblock Public Request Notice Selections for Source-Only Releases: Choosing Canonical Notice Families from Within-Release Status and Cross-Release Carryforward Tokens.
\newblock Series synthesis note (Synthesis~63), 2026. (This archive.)

\bibitem{series:publicrequestresponsemenus}
Anonymous.
\newblock Public Request Response Menus for Source-Only Releases: Smallest-Sufficient Follow-Up Handoffs for Basis Tokens, Visible Notices, Canonical Families, Family Choice, Scope, and Completion.
\newblock Series synthesis note (Synthesis~64), 2026. (This archive.)

\bibitem{series:requestfulfillmentcertificates}
Anonymous.
\newblock Request Fulfillment Certificates for Source-Only Releases: Exact Completion Criteria for Public Request Contracts and Typed Omitted-Support Cuts.
\newblock Series synthesis note (Synthesis~58), 2026. (This archive.)



\bibitem{series:publicrequestresponsepackets}
Anonymous.
\newblock Public Request Response Packets for Source-Only Releases: Exact Bounded Handoff Slices after Notice-Keyed Follow-Up Choice.
\newblock Series synthesis note (Synthesis~65), 2026. (This archive.)

\bibitem{series:publicrequestresponsepacketcarryforward}
Anonymous.
\newblock Public Request Response Packet Carryforward for Source-Only Successor Releases: Reuse Rules for Exact Notice-Keyed Handoff Slices under Release Evolution.
\newblock Series synthesis note (Synthesis~66), 2026. (This archive.)

\bibitem{series:publicrequestresponsepacketdeltas}
Anonymous.
\newblock Public Request Response Packet Delta Ledgers for Source-Only Successor Releases: Exact Refresh Instructions for Reused Notice-Keyed Handoff Slices.
\newblock Series synthesis note (Synthesis~67), 2026. (This archive.)

\bibitem{series:publicrequestresponsepacketrefreshnotices}
Anonymous.
\newblock Public Request Response Packet Refresh Notices for Source-Only Successor Releases: Exact Visible Reissue Sentences for Reused Notice-Keyed Handoff Slices.
\newblock Series synthesis note (Synthesis~68), 2026. (This archive.)

\bibitem{series:publicrequestresponsepacketrefreshnoticenormalforms}
Anonymous.
\newblock Public Request Response Packet Refresh Notice Normal Forms for Source-Only Successor Releases: Canonical Visible Reissue Families for Reused Notice-Keyed Handoff Slices.
\newblock Series synthesis note (Synthesis~69), 2026. (This archive.)

\bibitem{series:publicrequestresponsepacketrefreshnoticeselections}
Anonymous.
\newblock Public Request Response Packet Refresh Notice Selections for Source-Only Successor Releases: Choosing Canonical Visible Reissue Families from Packet-Local Delta Tokens.
\newblock Series synthesis note (Synthesis~70), 2026. (This archive.)

\bibitem{series:publicrequestresponsepacketrefreshresponsemenus}
Anonymous.
\newblock Public Request Response Packet Refresh Response Menus for Source-Only Successor Releases: Smallest-Sufficient Follow-Up Handoffs for Refreshed Visible Sentences.
\newblock Series synthesis note (Synthesis~71), 2026. (This archive.)


\bibitem{series:publicrequestresponsepacketrefreshresponsepackets}
Anonymous.
\newblock Public Request Response Packet Refresh Response Packets for Source-Only Successor Releases: Exact Bounded Handoff Slices after Concrete Follow-Up Choice about Visible Reissue Sentences.
\newblock Series synthesis note (Synthesis~72), 2026. (This archive.)

\bibitem{series:publicrequestresponsepacketrefreshresponsepacketclosureverdicts}
Anonymous.
\newblock Public Request Response Packet Refresh Response Packet Closure Verdicts for Source-Only Successor Releases: Certifying Stable Finite Basis after Exact Successor-Facing Handoff Slice Choice.
\newblock Series synthesis note (Synthesis~73), 2026. (This archive.)

\bibitem{series:pairedterminalcitationdiscipline}
Anonymous.
\newblock Paired Terminal Citation Discipline for Stable Review Spines: One Answer-Side Stop and One Request-Side Capstone for Terse Papers.
\newblock Series synthesis note (Synthesis~74), 2026. (This archive.)

\bibitem{series:pairedterminalcutoverrules}
Anonymous.
\newblock Paired Terminal Cutover Rules for Stable Review Spines: When Later Terse Papers Must Stop, Widen, or Reopen.
\newblock Series synthesis note (Synthesis~75), 2026. (This archive.)

\bibitem{series:theoremtopublicationspine}
Anonymous.
\newblock Stable Theorem-to-Publication Spines for Terse Anonymity Papers: Mathematical Roots, Worked Cuts, Release Evolution, Claim Objects, and Lineaged Releases.
\newblock Series synthesis note (Synthesis~76), 2026. (This archive.)




% Added in rev0806 citation-closure hygiene pass.
\bibitem{series:advantagecontracts}
Anonymous.
\newblock Advantage Contracts and Stealth Audits for Anonymous-DHT Privacy Claims.
\newblock Evaluation series Paper~2, 2026. (This archive.)

\bibitem{series:bossfightBaddendum}
Anonymous.
\newblock AnonDHT Dial Sheet: Supporting Evidence Tables and Pattern Map (Addendum).
\newblock Companion note (Boss Fight~B-A), 2026. (This archive.)

\bibitem{series:certifiedB}
Anonymous.
\newblock Routing-Signature Compression for Many-Testing-Safe Anonymous DHT Tuning.
\newblock Certified series Paper~B, The-One-True-Archive (2026).

\bibitem{series:certifiedmenus}
Anonymous.
\newblock Certified Menus for Anonymous DHT Lookups: Distribution-free parameter selection with simultaneous guarantees under many-testing.
\newblock Certified series Paper~A, 2026. (This archive.)

\bibitem{series:congestion1}
Anonymous.
\newblock Congestion-EQ: Total-Variation Contracts for Queue/RTT Leakage in Anonymous Overlays and Anonymous DHT Lookups.
\newblock Congestion series Paper~1, The-One-True-Archive (2026).

\bibitem{series:eval1}
Anonymous.
\newblock Notarized Anonymity Certificates: Attempt Logs, Spending Discipline, and Verifier Rules.
\newblock Evaluation series Paper~1, 2026. (This archive.)

\bibitem{series:opB}
Anonymous.
\newblock Auditable Trace Privacy for Anonymous DHTs: Runtime Conformance, Plan Pinning, and Epoch Receipts.
\newblock Operational series Paper~B, 2026. (This archive.)

\bibitem{series:operationalA}
Anonymous.
\newblock Schedule-Only Retries in Anonymous DHTs: When Retry Logic Becomes a Linkability Channel.
\newblock Operational series Paper~A, The-One-True-Archive (2026).

\bibitem{series:operationalB}
Anonymous.
\newblock Auditable Trace Privacy with Abort: Plan-Pinning and Drift Receipts for Side-Channel Resistant Transport.
\newblock Operational series Paper~B, The-One-True-Archive (2026).

\bibitem{series:pscq}
Anonymous.
\newblock \emph{PSC-Q: Public Service Calendars and Slack-Harvesting Substitution Padding for Congestion-Private Anonymous Overlays}.
\newblock Companion paper in this unified archive.

\bibitem{series:recert}
Anonymous.
\newblock Disagreement-Gated Recertification and Boundary Audits for Anonymous DHT Routers.
\newblock Certified series Paper~C, 2026. (This archive.)

\bibitem{series:release}
Anonymous.
\newblock Release Property, Ledger Receipts, and Public Lineage for Anonymity Claims.
\newblock Release-and-Destination series Paper~A, The-One-True-Archive (2026).

\bibitem{series:releaseA}
Privacy as a Release Property: Proof-Carrying Anonymity Receipts for Anonymous DHTs.
\newblock Release and Destination series Paper~A, 2026. (This archive.)

\bibitem{series:replayrecipe}
Anonymous.
\newblock Auditor Replay for Anonymous-DHT Receipt Line Items: A Mechanical Checklist for Digest-Bound Claims.
\newblock Series synthesis note (Synthesis~20), 2026. (This archive.)

\bibitem{series:schedretries}
Anonymous.
\newblock \emph{Schedule-Only Retries for Anonymous DHT Lookups}.
\newblock Operational series Paper~A, The-One-True-Archive (2026).

\bibitem{series:wcongestion}
Anonymous.
\newblock W-Congestion-EQ: Verifiable Congestion Budgets for Anonymous Overlays.
\newblock Congestion series Paper~3, The-One-True-Archive (2026).

\bibitem{series:workedexample}
Anonymous.
\newblock Worked Example: Instantiating a Tiered Reader-Privacy Receipt.
\newblock Synthesis~17, The-One-True-Archive (2026).

\bibitem{state1}
Anonymous.
\newblock State-Dependent Anonymity in Anonymous DHTs: Witnesses, Amplification, and Stability.
\newblock 2026.
\newblock AnonDHT State series Paper~1 (this archive).

\bibitem{state2}
Anonymous.
\newblock Marginally Uniform Committee Contact in Anonymous DHT Lookups: Availability Floors, Perfect-Leakage Counterexamples, and Abort-Safe Equalization.
\newblock 2026.
\newblock AnonDHT State series Paper~2 (this archive).

\bibitem{state3}
Anonymous.
\newblock Closed-View Auditing for Stateful Anonymity: Odometers, Alarms, and Machine-Checkable CPPCs.
\newblock 2026.
\newblock AnonDHT State series Paper~3 (this archive).

\end{thebibliography}

\end{document}

